% concatenated from example_paper.tex, intro.tex
\documentclass{article}

\usepackage{microtype}
\usepackage{graphicx}
\usepackage{subcaption} % Modern replacement for subfigure
\usepackage{booktabs}   % for professional tables
\usepackage{hyperref}   % hyperlinks

\newcommand{\theHalgorithm}{\arabic{algorithm}}

\usepackage[accepted]{icml2026}

\usepackage{amsmath}
\usepackage{amssymb}
\usepackage{mathtools}
\usepackage{amsthm}
\usepackage{amsfonts}
\usepackage{bm}            % Bold math
\usepackage{nicefrac}      % Compact fractions
\usepackage{xcolor}        % Colors
\usepackage{url}

% Tables
\usepackage{multirow}
\usepackage{makecell}

\usepackage{threeparttable}
\usepackage{array}
\usepackage{tabularx}

% Misc
\usepackage{xspace}
\usepackage{enumerate}
\usepackage[inline]{enumitem} 
\usepackage{bbold}
\usepackage{epstopdf}
\usepackage{rotating}
\usepackage{mathrsfs}
\usepackage[normalem]{ulem} % underlining
\usepackage{physics}        % Handy physics notation

% ------------------------------------------------------------------
% 3. Algorithm Setup (CRITICAL PART)
% ------------------------------------------------------------------
\usepackage{algorithm}
\usepackage{algorithmic}

% ------------------------------------------------------------------
% 4. Custom Macros (Definitions)
% ------------------------------------------------------------------
\allowdisplaybreaks

% Box Utilities
\newcommand*{\colorboxedAux}[3]{%
  \begingroup
    \colorlet{cb@saved}{.}%
    \color#1{#2}%
    \boxed{%
      \color{cb@saved}%
      #3%
    }%
  \endgroup
}
\newcommand*{\colorboxed}{}
\def\colorboxed#1#{%
  \colorboxedAux{#1}%
}

% --- Algorithm Commands Adjustments ---
\renewcommand{\algorithmicrequire}{\textbf{Input:}}
\renewcommand{\algorithmicensure}{\textbf{Initial values:}}
\newcommand{\rgrad}{\mathrm{grad}}
\newcommand{\Log}{\mathrm{Log}}
\newcommand{\pert}{\nu}
\DeclareMathOperator*{\argmin}{argmin}
\DeclareMathOperator{\sign}{sign}
\DeclareMathOperator{\dom}{dom}

% ------------------------------------------------------------------
% 5. Theorem & Cleveref (Order Matters: Cleveref Last)
% ------------------------------------------------------------------
\usepackage[capitalize,noabbrev]{cleveref}

\theoremstyle{plain}
\newtheorem{theorem}{Theorem}[section]
\newtheorem{proposition}[theorem]{Proposition}
\newtheorem{lemma}[theorem]{Lemma}
\newtheorem{corollary}[theorem]{Corollary}
\theoremstyle{definition}
\newtheorem{definition}[theorem]{Definition}
\newtheorem{assumption}[theorem]{Assumption}
\theoremstyle{remark}
\newtheorem{remark}[theorem]{Remark}

\usepackage[textsize=tiny]{todonotes}

\icmltitlerunning{Riemannian Dueling Optimization}

\begin{document}

\twocolumn[
  \icmltitle{Riemannian Dueling Optimization}

  \icmlsetsymbol{equal}{*}

  \begin{icmlauthorlist}
    \icmlauthor{Yuxuan Ren}{rice}
    \icmlauthor{Abhishek Roy}{tamu}
    \icmlauthor{Shiqian Ma}{rice}
  \end{icmlauthorlist}

  \icmlaffiliation{rice}{Department of Computational Applied Mathematics and Operations Research, Rice University, Houston, TX, USA}
  \icmlaffiliation{tamu}{Department of Statistics, Texas A\&M University, College Station, TX, USA}

  \icmlcorrespondingauthor{Yuxuan Ren}{yr21@rice.edu}
  \icmlcorrespondingauthor{Abhishek Roy}{abhishekroy@tamu.edu}
  \icmlcorrespondingauthor{Shiqian Ma}{sqma@rice.edu}

  \icmlkeywords{Machine Learning, ICML}

  \vskip 0.3in
]


\printAffiliationsAndNotice{} 


\begin{abstract}
Dueling optimization considers optimizing an objective with access to only a comparison oracle of the objective function. It finds important applications in emerging fields such as recommendation systems and robotics. Existing works on dueling optimization mainly focused on unconstrained problems in the Euclidean space. In this work, we study dueling optimization over Riemannian manifolds, which covers important applications that cannot be solved by existing dueling optimization algorithms. In particular, we propose a Riemannian Dueling Normalized Gradient Descent (RDNGD) method and establish its iteration complexity when the objective function is geodesically $L$-smooth or geodesically (strongly) convex. We also propose a projection-free algorithm, named Riemannian Dueling Frank–Wolfe (RDFW) method, to deal with the situation where projection is prohibited. We establish the iteration and oracle complexities for RDFW. We illustrate the effectiveness of the proposed algorithms through numerical experiments on both synthetic and real applications. 
\end{abstract}

\section{Introduction}
Many modern learning tasks involve settings where gradients or even function values are inaccessible, 
and the only feasible feedback comes in the form of pairwise preferences or comparison -- often referred to as \emph{dueling feedback}. 
For example, users in recommendation systems typically express relative judgments such as 
``item A is preferred over item B''; in robotics, human supervisors often compare trajectories rather than assign scalar rewards \citep{yang2025trajectory,jain2013learning}; 
and in representation learning, a downstream classifier may indicate that one projection preserves class structure better than another \citep{ventocilla2020comparative,morariu2021dumbledr}. 
Importantly, in many of these applications the decision space itself is inherently non-Euclidean: recommendation models often rely on hierarchical embeddings in hyperbolic space \citep{chamberlain2019scalable,shimizu2024fashion}, 
trajectory optimization in robotics often requires optimizing over special orthogonal groups $\mathrm{SO}(3)$ \citep{watterson2018trajectory}, and projection matrices for representation learning are naturally constrained to the Stiefel manifold \citep{boumal2023intro}. In addition, solutions to high-dimensional learning problems frequently lie on low-dimensional manifolds \citep{garipov2018loss,hu2022lora}. Furthermore, several preference elicitation frameworks require constrained optimization where feasible solutions belong to constraint sets such as sphere, and simplex. These problems can also be reformulated as optimization over a manifold \citep{boumal2023intro}. Classical problems such as the Karcher mean problem where the feedback is provided by human and they are not aware of the loss can also be formulated as a Riemannian dueling optimization problem. Motivated by these applications, we study the Riemannian dueling optimization problem in the following form:
\begin{equation} 
\min_{x\in\mathcal{X}\subseteq\mathcal{M}} \quad f(x), 
\label{eq:target1.1}
\end{equation}
where $\mathcal{M}$ is a Riemannian manifold with dimension $d$, and $f:\mathcal{X}\subseteq\mathcal{M}\to\mathbb{R}\cup\{\infty\}$ is a smooth function. Let $x^*$ denote the minimizer of \eqref{eq:target1.1}, and $f^* := f(x^*)$. We assume that we have access to only a pairwise comparison oracle $\mathcal{Q}_f(x,y)$ between two queried points $x,y\in\mathcal{M}$ given by 
\begin{align}\label{eq:oracle}
    \mathcal{Q}_f(x,y)=2*\mathbb{1}(f(x)>f(y))-1,
\end{align}
where $\mathbb{1}(\cdot)$ is the indicator function, which equals 1 if the statement is true and 0 otherwise. Note that we do not have access to the function values $f(x)$ and $f(y)$.

\subsection{Review of Euclidean Dueling Optimization}
%The main challenge posed by dueling feedback setting is that it destroys the gradient magnitude information unlike first-order, and zeroth-order setting. Even within preference-based optimization, one can sometimes recover gradient magnitude information, e.g., under the popular Bradley--Terry model the probability of observing a specific preference outcome depends on the strength of the utility difference between two queried points, thereby allowing estimation of magnitude of utility gradient to some extent \cite{saha2022generalpreference,wang2023rlhf}. In contrast, under the feedback model \eqref{eq:oracle}, only the sign of the comparison is revealed, and all information about gradient magnitude is lost. However, the gradient direction can still be estimated. 

%While several algorithms have been developed for this setting in Euclidean space, their extension to Riemannian manifolds reveals intrinsic limitations and computational bottlenecks:

In this subsection, we briefly review the literature on Euclidean dueling optimization. \cite{lobanov2024acceleration} and \cite{chervonenkis2024nesterov} proposed coordinate descent algorithms for dueling optimization. \cite{bergou2020stochastic} introduced a stochastic three points method which is a direct search algorithm based on comparison oracles. Variants of this method were further studied in \citep{boucherouite2024minibatch,kadi2025almost,tkachenkothree}. These methods typically rely on identifying the best candidate among the current iterate and its neighbors, which requires multiple calls to the comparison oracle in each iteration. The main difficulty in dueling optimization is how to efficiently estimate the gradient direction using comparison oracle, and recent advances have explored sophisticated techniques for explicit gradient approximation. \cite{Cai_2022} proposed one-bit and comparison-based gradient estimators; \cite{tang2024zerothorder} proposed to estimate gradient using ranking oracles; \cite{zhang2024comparisons} proposed a binary search method to compute the gradient estimator using comparison oracle. In our work, we extend the two-point random perturbation approach for gradient estimation proposed by \citet{pmlr-v139-saha21b} to the Riemannian setting, as it requires only a single random direction and one comparison per iteration. However, extending this framework is far from a trivial generalization, as the presence of curvature breaks standard Euclidean trigonometry and linearization.

%\begin{itemize}[leftmargin=*]
%    \item \textbf{Coordinate Descent and Line-Search:} A random coordinate descent algorithm based on line-search was proposed in \cite{lobanov2024acceleration} that achieves better dimension dependence compared to the former ones. This approach requires computing the minimum of the function along a fixed direction that requires multiple function evaluations. \cite{chervonenkis2024nesterov} demonstrated that acceleration is possible for coordinate descent in a restricted setting, and relies on line search. This is not a problem in Euclidean setting but requires multiple computationally expensive exponential map computations in the Riemannaian setting. 

    
    
 %   \item \textbf{Direct Search and Reduction Frameworks:} A line of research explores derivative-free optimization via direct search and comparison-based strategies. The Stochastic Three Points (STP) method \cite{bergou2020stochastic}, along with its variants incorporating minibatching \cite{boucherouite2024minibatch}, theoretical guarantees \cite{kadi2025almost}, and adaptive sampling \cite{tkachenkothree}, operates as a sign-based direct search algorithm. These methods typically rely on identifying the best candidate among the current iterate and its neighbors, necessitating multiple calls to the comparison oracle per iteration to filter random directions. Similarly, general reduction frameworks like \cite{slavin2022adapting} adapt standard zeroth-order algorithms to the comparison-based setting by simulating function value rankings. However, this strategy is costly because it requires many pairwise comparisons to simulate a single function evaluation. 
    
    
  %  \item \textbf{Gradient Estimators:} Recent advances have explored sophisticated techniques for explicit gradient approximation, though often at prohibitive costs. \cite{Cai_2022} proposed estimators based on 1-bit compressed sensing in the strongly convex setting and \cite{tang2024zerothorder} proposed to estimate gradient using ranking oracles. Despite their methodological differences, both approaches inherently rely on aggregating information from a large number of samples per iteration. In the Riemannian setting, this high sample complexity either leads to massive evaluations of the computationally expensive exponential map to generate candidates, or requires an excessive number of oracle calls. \cite{zhang2024comparisons} utilize an orthonormal basis to control variance, but their method is computationally expensive due to the need for per-sample dueling and basis maintenance. This results in heavy oracle complexity and a severe dependency on the problem dimension, which is impractical for high-dimensional tasks. The two-point random perturbation approach proposed by \citet{pmlr-v139-saha21b} is particularly suitable for the Riemannian setting due to its minimal oracle overhead. Requiring only a single random direction and comparison per iteration is essential for manifolds, where repeated exponential map evaluations are computationally demanding. However, extending this framework is far from a trivial generalization.The presence of curvature breaks standard Euclidean trigonometry and linearization.  we bridge this gap by generalizing dueling optimization to the Riemannian manifold setting. We establish a comprehensive theoretical framework for Riemannian dueling optimization, encompassing both unconstrained and constrained scenarios. 
%\end{itemize}

% In this paper, we bridge this gap by generalizing dueling optimization to the Riemannian manifold setting. We establish a comprehensive theoretical framework for Riemannian dueling optimization, encompassing both unconstrained and constrained scenarios. Specifically, we offer projected-type methods for cases where optimal convergence rates are prioritized, and a novel projection-free framework to handle manifold constraints where the standard projection operator is computationally intractable.

% we consider the setting where the learner has access only to pairwise comparison feedback: for any two queried points $x,y\in\mathcal{M}$, the oracle reveals whether $f(x)>f(y)$.  In contrast, prior works on Riemannian optimization assume access to either first-order feedback consisting of gradients \cite{zhang2016riemannian,sato2019riemannian,alimisis2021momentum,zhou2025inexact,he2025relatively,sahinoglu2025decentralized,wang2025federated,zhang2018r,pmlr-v125-ahn20a,weber2022projection}, or zeroth-order feedback, consisting of function evaluations \cite{wang2021greene,li2023stochastic,wang2023sharp,li2024zeroth,zhou2024inexact,he2024riemannian}. In contrast, our setting allows access to only pairwise comparison $\cP(x,y)$ between any two queried points $x,y\in\mathcal{M}$ given by 
% \begin{align}\label{eq:oracle}
%     \cP(x,y)=\mathbb{1}(f(x)>f(y)),
% \end{align}
% where $\mathbb{1}(\cdot)$ is the indicator function.

% The main challenge posed by dueling feedback setting is that it destroys the gradient magnitude information unlike first-order, and zeroth-order setting. Even within preference-based optimization, one can sometimes recover gradient magnitude information, e.g., under the popular Bradley--Terry model the probability of observing a specific preference outcome depends on the strength of the utility difference between two queried points, thereby allowing estimation of magnitude of utility gradient to some extent \cite{saha2022generalpreference,wang2023rlhf}. In contrast, under the feedback model \eqref{eq:oracle}, only the sign of the comparison is revealed, and all information about gradient magnitude is lost. However, the gradient direction can still be estimated. Several gradient direction estimator-based algorithms have been proposed recently for unconstrained optimization with dueling feedback in the Euclidean space \citep{pmlr-v139-saha21b,blum2024duelingmonotone,zhang2024comparisons,zhang2024randomized}. A random coordinate descent algorithm based on line-search was proposed in \cite{lobanov2024acceleration} that achieves better dimension dependence compared to the former ones. This approach requires computing the minimum of the function along a fixed direction that requires multiple function evaluations. \cite{chervonenkis2024nesterov} demonstrated that acceleration is possible for coordinate descent in a restricted setting, and relies on line search. This is not a problem in Euclidean setting but requires multiple computationally expensive exponential map computations in the Riemannaian setting. However, none of these works addresses optimization on Riemannian manifold. Moreover, all the algorithms proposed in Euclidean settings are for unconstrained optimization. 

% In this paper, we take a step towards bridging this gap by establishing theoretical guarantees for dueling feedback optimization in both unconstrained and constrained settings. Our contributions encompass the general unconstrained case, while for constrained problems, we offer projected-type methods for faster convergence and projection-free algorithms to handle complex constraints where projection is intractable.

% \paragraph{Technical Challenges} The first technical challenge in designing optimization algorithms with dueling feedback on manifolds is the absence of theoretically sound methods for estimating Riemannian gradient directions. We address this by proposing a two-point perturbation estimator, in the spirit of \citet{pmlr-v139-saha21b} for the Euclidean case. Unlike their setting, the manifold structure requires care since perturbed points need not remain on the manifold. We establish a sharper bias bound for our estimator (Lemma~\ref{lem:grad_with_bias}), improving upon \citet{pmlr-v139-saha21b}; our analysis also yields tighter bounds in Euclidean space. 

% For the projected normalized gradient method with dueling oracle, we analyze both constant stepsize and a diminishing cosine-annealing stepsize; With \Cref{lem:cosine_annealing_stepsizes} to handle the cosine schedule, the proof technique is essentially identical and the convergence rate differs only by a constant factor. The motivation for the diminishing stepsize is practical: constant stepsizes perform well on smaller-dimensional tests, and cosine annealing stepsize tends to work better on larger instances. The analysis for the cosine annealing stepsize also provides more flexibility for parameter tuning.

% We introduce a Riemannian Frank–Wolfe variant that employs an increasing batch-size strategy to progressively reduce variance. To the best of our knowledge, this is the first projection-free algorithm with dueling feedback for constrained optimization, a gap that remained open even in the Euclidean setting. In this case, controlling both the bias and the variance of the gradient direction estimator is essential. This strategy is particularly effective as it ensures that the iteration complexity is independent of the dimension. Furthermore, our analysis directly tackles the geometric challenges imposed by curvature. By addressing the problem directly on the manifold, our complexity bounds depend only on the intrinsic dimension of the manifold rather than the significantly larger dimension of the ambient space.



\subsection{Motivating Examples}\label{sec:example}
% of Riemannian Dueling Optimization}

We now discuss two motivating examples for Riemannian dueling optimization. 

{\bf Attack on Deep Neural Network.}
Attacks on deep neural networks (DNNs) aim to introduce small, human-imperceptible perturbations to an input so that the model produces an incorrect label. In realistic black-box scenarios, the attacker typically has only limited query access to the model and therefore cannot rely on gradient information. Moreover, the attacker may not even have access to a clean or reliable loss function. For example, prediction APIs often return only labels or heavily processed scores \citep{xie2018, s.2018stochastic, athalye2018obfuscatedgradientsfalsesense}, so function values are either unavailable or not reliable. In such cases, it is more natural to ask which of two perturbed images is more adversarial, and to use these pairwise preferences to drive the optimization. % instead of relying on explicit loss values. 
Furthermore, natural images are widely believed to concentrate near low-dimensional manifolds \citep{ws04}. %, while most existing zeroth-order attacks operate directly in the ambient Euclidean space and produce perturbed images that need not remain within the natural data domain. Such perturbations that deviate from the manifold structure often lead to unrealistic samples and inefficient search dynamics. 
Similar to the optimal $\ell_2$-norm attack framework established in \citep{lyu2015unifiedgradientregularizationfamily} and further explored in \citep{li2023stochastic}, we model the constraint manifold $\mathcal{M}$ as a sphere for simplicity. %Formally, the perturbation set is defined as $\mathcal{M} = \mathcal{S}_d(R):=\{x\in \mathbb{R}^d:\|x\|_2=R\}$ where $R$ represents the sphere's radius. This formulation leads to the following constrained optimization problem:
This leads to the following dueling optimization in the form of \eqref{eq:target1.1}:
\begin{equation*}
    \max_{x \in \mathcal{M}} \quad f(x) := \mathcal{L}_{w}(x_0 + x),
\end{equation*}
where $\mathcal{L}_w$ denotes the loss function of a DNN with weights $w$, $x_0$ is the image to attack, and the decision variable $x$ represents the adversarial perturbation. %A primary advantage of this spherical constraint is that it naturally restricts the perturbed image to a bounded feasible domain, ensuring that the perturbation maintains a constant and controlled $\ell_2$ distance from the original input.


%To address this limitation, we develop a Riemannian dueling attack that performs adversarial optimization on an explicit manifold representation of the data. 

% Distinct from existing Riemannian zeroth-order approach \cite{li2023stochastic}, our method leverages a dueling comparison oracle to infer descent directions on the manifold, enabling efficient black-box attacks while preserving the intrinsic geometry of natural images.


{\bf Horizon Leveling.}
Horizon leveling is a fundamental task in computer vision, where the goal is to estimate and correct the image horizon. It plays a key role in applications such as single-image camera calibration \citep{HoldGeoffroy2018} and visual navigation \citep{Ettinger2002}. 
%Recent deep models for horizon line estimation have been shown to work reliably in the wild, even under complex natural scenes and cluttered geometries \citep{Workman2016}. Given an input image, we model corrections as rotations on the group $\mathrm{SO}(2)$ and seek an angle that makes the scene ``look level.'' 
Formally, the horizon tilt in an image can be represented by a rotation matrix $R_{\text{tilt}} \in \mathrm{SO}(2)$, where a perfectly level horizon corresponds to the identity matrix $I$. %Here, $\mathrm{SO}(2) = \{ R \in \mathbb{R}^{2 \times 2} \mid R^\top R = I, \det(R) = 1 \}$ denotes the Special Orthogonal group of 2D rotations. 
%\sout{Our objective is to determine a correction rotation $R \in \mathrm{SO}(2)$ that aligns the tilted image with the horizon. We quantify the misalignment of a candidate correction using the cost function of $f(R) := \| R R_{\text{tilt}} - I \|_F^2$.}
In many realistic settings, however, an explicit loss function or ground-truth labels may be unavailable. Instead, we observe only pairwise comparisons from the loss function $f(R)$ indicating which of two rotated versions appears more level. 
%Instead, we only observe comparative feedback: for example, a human or a black-box system can say which of two rotated versions appears more horizontally aligned, but does not provide gradients or calibrated scores. 
%This naturally leads to the following Riemannian dueling optimization: %a Riemannian optimization problem on the $\mathrm{SO}(2)$ manifold:
%follows which is in the form of \eqref{eq:target1.1}:
This leads naturally to the following Riemannian dueling optimization problem over $\mathrm{SO}(2)$, where the goal is to find the correction $R^\star$ that best compensates for the observed tilt $R_{\text{tilt}}$.
\begin{equation*}
    \min_{R \in \mathrm{SO}(2)} f(R).
\end{equation*}
%The goal is to find the correction $R^\star \in \mathrm{SO}(2)$ that minimizes the misalignment cost, thereby best compensating for the observed tilt $R_{\text{tilt}}$. 
\begin{remark}
Our applications illustrate distinct feedback sources: machine-generated adversariality (DNN attacks) versus human preferences (horizon leveling). This demonstrates that our framework is agnostic to the comparison source, unifying both model-based and human-based objectives within a single scheme.
\end{remark}



\subsection{Contributions}
%This paper is the first one to consider Riemannian dueling optimization, where only a comparison oracle of objective function is available. It thus solves Riemannian optimization problems that cannot be handled by classical Riemannian optimization algorithms. The main contributions of this paper are summarized as follows.
In this paper, we introduce Riemannian dueling optimization, where only a comparison oracle of the objective function is available, and thus classical Riemannian optimization algorithms fail. Our main contributions are as follows.
\begin{itemize}[leftmargin=*]
    \item We propose the first Riemannian dueling normalized gradient descent method (RDNGD) for Riemannian dueling optimization and provide its iteration and oracle complexities for geodesically $L$-smooth objectives, as well as for geodesically $L$-smooth and (strongly) geodesically convex objectives.
%    We solve the problems of Riemannian geodesically convex and nonconvex optimization with a comparison oracle of single-point feedback. To the best of our knowledge, this is the first algorithm that provides convergence guarantees for comparison-oracle–based optimization on Riemannian manifolds, in both convex and nonconvex settings. 
%    \item We develop a novel analysis with an improved bound on the bias in the gradient direction estimation based on comparison oracles. This analysis can be applied to both the Euclidean and Riemannian cases. {\color{red}remove}%Although the improvement does not change the order of the overall complexity, it refines the constant factors and offers more flexible strategy for selecting other algorithmic parameters.
%    \item We propose two Riemannian normalized gradient descent methods with dueling oracle, which utilize projection operators to maintain feasibility within the manifold constraints. We provide complexity result for the different function classes of (i) geodesically $\beta$-smoothness (ii) geodesically $\beta$-smoothness and geodesically convexity (iii) geodesically $\beta$-smoothness and geodesically $\alpha$-strongly convexity. The complexity result of each case matches its best Euclidean counterpart.
    % \item We propose one algorithm in the Riemannian Frank–Wolfe family for geodesically convex optimization using a dueling oracle. The algorithm leverages a batch-averaged search direction and a variance-reduction hybrid direction, and we derive its convergence guarantees. To our knowledge, this is the first time for the Frank–Wolfe method with dueling oracle to be explored, even in the Euclidean space.   
    \item We propose a Riemannian dueling Frank-Wolfe method (RDFW), which is projection-free, to handle cases where projection is computationally prohibitive. %The algorithm integrates a batch-averaged search direction with a variance-reduction hybrid strategy to effectively optimize objectives using only preference-based feedback. 
    We establish iteration and oracle complexities of RDFW in the geodesically convex setting, which is the first convergence result for projection-free dueling optimization over manifolds.
\end{itemize}
Overall, this work bridges two active areas, preference-based optimization and Riemannian optimization, by showing how optimization can proceed reliably on manifolds when only comparison oracles are available. Our main results are summarized in Table~\ref{tab:cases}. 
%RDNGD, RRDNGD and RDFW are our algorithms for Riemannian dueling optimization for different problem classes, and $\beta$-NGD and ($\alpha$,$\beta$)-NGD are existing algorithms from \citep{pmlr-v139-saha21b} that are designed for Euclidean dueling optimization. 
Our \texttt{RDNGD} and \texttt{RRDNGD} algorithms can be viewed as Riemannian generalizations of the NGD methods proposed in \citep{pmlr-v139-saha21b}. While we employ a similar gradient direction estimator, our contribution extends significantly beyond a trivial generalization to the Riemannian setting. In fact, when reducing to the Euclidean setting, we also provide better results comparing with the results in \citep{pmlr-v139-saha21b}, and we will specify them when we present the specific results. %More specifically, (i) We provide a sharper analysis of the gradient estimator by optimizing the choice of the perturbation parameter $\nu$ and improving theoretical constants; (ii) \textbf{Sharper Convergence Bounds:} For the strongly g-convex setting using the restart strategy ({RRDNGD}), our modified parameter selection yields a superior convergence rate. Specifically, we improve the complexity's dependence on the initial squared distance $D$ from linear ($\mathcal{O}(D)$) to logarithmic ($\mathcal{O}(\log D)$). %(iii) \textbf{Extension to Projection-Free Optimization:} We further demonstrate the versatility of our estimator by integrating it into the Frank-Wolfe framework. This results in the Riemannian Dueling Frank-Wolfe (RDFW) algorithm, which efficiently solves constrained g-convex problems without requiring computationally expensive projections.

%\subsection{Problem Settings}
%We summarize the problem settings addressed in this work. All algorithms considered in later sections aim to solve target problem~\eqref{eq:target1.1}
% \begin{equation*}
% \min_{x\in\mathcal{X}\subset\mathcal{M}} \quad f(x),
% % \label{eq:general_problem_setting} 
% \end{equation*}
% but differ in the assumptions imposed on the objective function and feasible set.

%Table~\ref{tab:cases} summarizes the four problem classes we study. These include (1) unconstrained nonconvex optimization, (2) convex constrained convex optimization, (3) strongly convex constrained convex optimization, and (4) convex constrained convex optimization with expensive projections, together with the assumptions required by each algorithm.

%The table is divided into two distinct groups: the methods proposed in this work (RDNGD, RRDNGD, RDFW) and the relevant existing literature ($\beta$-NGD, $(\alpha,\beta)$-NGD), which provides a basis for comparison. Although RDNGD appears twice in the table (representing nonconvex and convex settings), the underlying algorithm is identical in both cases; the only distinction lies in whether the projection operator is required. In the unconstrained setting, the projection reduces to the identity map, and thus no constraint handling is needed.



% \begin{table*}[t]
% \centering
% \caption{Comparison of the proposed algorithms with existing methods.}
% \label{tab:cases}
% \resizebox{\textwidth}{!}{%
%     \renewcommand{\arraystretch}{1.35}
%     \begin{tabular}{c|c|c|c|c|c} 
%     \hline
%     \textbf{Algorithm} &
%     \textbf{\makecell{Requirements \\ on $f$}} & 
%     \textbf{\makecell{Requirements \\ on $X$}} & 
%     \textbf{\makecell{Constraint \\ Handling}} &
%     \textbf{\makecell{Iteration \\ Complexity}} &
%     \textbf{\makecell{Oracle \\ Complexity}} \\
%     \hline

%     % --- Proposed Methods ---
%     \textbf{RDNGD (nonconvex)} &
%     $L$-smooth &
%     Unconstrained ($\mathcal{X} = \mathcal{M}$) &
%     / &
%     $\mathcal{O}(d\epsilon^{-2})$ &
%     $\mathcal{O}(d\epsilon^{-2})$
%     \\
%     \hline

%     \textbf{RDNGD (convex)} &
%     \begin{tabular}{c}
%     $L$-smooth \\
%     g-convex
%     \end{tabular} &
%     Geodesically uniquely convex &
%     Projection &
%     $\mathcal{O}(d\epsilon^{-1})$ &
%     $\mathcal{O}(d\epsilon^{-1})$
%     \\
%     \hline

%     \textbf{RRDNGD} &
%     \begin{tabular}{c}
%     $L$-smooth \\
%     strongly g-convex
%     \end{tabular} &
%     Geodesically uniquely convex &
%     Projection &
%     $\mathcal{O}(d\log(1/\epsilon))$ &
%     $\mathcal{O}(d\log(1/\epsilon))$
%     \\
%     \hline

%     \textbf{RDFW} &
%     \begin{tabular}{c}
%     $L$-smooth \\
%     g-convex
%     \end{tabular} &
%     Geodesically uniquely convex &
%     LMO &
%     $\mathcal{O}(\epsilon^{-1})$ &
%     $\mathcal{O}(d\epsilon^{-2})$
%     \\
%     \hline
%     \hline

%     % --- Existing Methods ---
%     \textbf{$\beta$-NGD }\citep{pmlr-v139-saha21b} &
%     \begin{tabular}{c}
%     $L$-smooth \\
%     convex
%     \end{tabular} &
%     Unconstrained &
%     / &
%     $\mathcal{O}(d\epsilon^{-1})$ &
%     $\mathcal{O}(d\epsilon^{-1})$
%     \\
%     \hline

%     \textbf{$(\alpha,\beta)$-NGD }\citep{pmlr-v139-saha21b} &
%     \begin{tabular}{c}
%     $L$-smooth \\
%     strongly convex
%     \end{tabular} &
%     Unconstrained &
%     / &
%     $\mathcal{O}(d\log(1/\epsilon))$ &
%     $\mathcal{O}(d\log(1/\epsilon))$
%     \\
%     \hline

%     \end{tabular}%
% }
% \end{table*}
\begin{table*}[t]
\centering
\caption{Comparison of the proposed algorithms with existing methods.}
\label{tab:cases}
\resizebox{\textwidth}{!}{%
    \renewcommand{\arraystretch}{1.35}
    \begin{tabular}{cccccc} 
    \toprule
    \textbf{Algorithm} &
    \textbf{\makecell{Requirements \\ on $f$}} & 
    \textbf{\makecell{Requirements \\ on $\mathcal{
    X}$}} & 
    \textbf{\makecell{Constraint \\ Handling}} &
    \textbf{\makecell{Iteration \\ Complexity}} &
    \textbf{\makecell{Oracle \\ Complexity}} \\
    \midrule

    % --- Proposed Methods ---
    \textbf{RDNGD (nonconvex)} &
    g-$L$-smooth &
    Unconstrained ($\mathcal{X} = \mathcal{M}$) &
    / &
    $\mathcal{O}(d\epsilon^{-2})$ &
    $\mathcal{O}(d\epsilon^{-2})$
    \\
    \midrule

    \textbf{RDNGD (convex)} &
    \begin{tabular}{c}
    g-$L$-smooth,
    g-convex
    \end{tabular} &
    Geodesically uniquely convex &
    Projection &
    $\mathcal{O}(d\epsilon^{-1})$ &
    $\mathcal{O}(d\epsilon^{-1})$
    \\
    \midrule

    \textbf{RRDNGD} &
    \begin{tabular}{c}
    g-$L$-smooth, 
    strongly g-convex
    \end{tabular} &
    Geodesically uniquely convex &
    Projection &
    $\mathcal{O}(d\log(1/\epsilon))$ &
    $\mathcal{O}(d\log(1/\epsilon))$
    \\
    \midrule

    \textbf{RDFW} &
    \begin{tabular}{c}
    g-$L$-smooth, 
    g-convex
    \end{tabular} &
    Geodesically uniquely convex &
    LMO &
    $\mathcal{O}(\epsilon^{-1})$ &
    $\mathcal{O}(d\epsilon^{-2})$
    \\
    \midrule
    \midrule

    % --- Existing Methods ---
    \textbf{$\beta$-NGD }\citep{pmlr-v139-saha21b} &
    \begin{tabular}{c}
    $L$-smooth,
    convex
    \end{tabular} &
    Unconstrained &
    / &
    $\mathcal{O}(d\epsilon^{-1})$ &
    $\mathcal{O}(d\epsilon^{-1})$
    \\
    \midrule

    \textbf{$(\alpha,\beta)$-NGD }\citep{pmlr-v139-saha21b} &
    \begin{tabular}{c}
    $L$-smooth,
    strongly convex
    \end{tabular} &
    Unconstrained &
    / &
    $\mathcal{O}(d\log(1/\epsilon))$ &
    $\mathcal{O}(d\log(1/\epsilon))$
    \\
    \bottomrule
    \end{tabular}%
}
\end{table*}


% \begin{remark}[Connections to Existing Riemannian Optimization Methods]
% While our primary focus is on the dueling feedback setting, it is instructive to situate our results within the broader context of existing Riemannian optimization methods:
% \begin{itemize}[leftmargin=*]
%     \item \textbf{Comparison with Existing Riemannian Zeroth-Order Methods:} To the best of our knowledge, there are currently no existing algorithms specifically designed for the dueling feedback setting on Riemannian manifolds. We therefore compare our \texttt{RDNGD} with \texttt{ZO-RGD} \cite{li2023stochastic}, a representative zeroth-order Riemannian method. Although both achieve an oracle complexity of $\mathcal{O}(d\epsilon^{-2})$, \texttt{ZO-RGD} requires full function value access. Our framework establishes that the same convergence rate is attainable using only $1$-bit sign information, thereby achieving comparable efficiency with a significantly weaker feedback signal.
    
%     \item \textbf{Comparison with Riemannian Frank-Wolfe Methods:} For constrained optimization, our proposed \texttt{RDFW} stands as the first projection-free algorithm under the dueling feedback setting. It achieves an iteration complexity of $\mathcal{O}(\epsilon^{-1})$, matching the first-order Riemannian Frank-Wolfe (\texttt{RFW}) \cite{weber2022projection}. The gap in oracle complexity ($\mathcal{O}(d\epsilon^{-2})$ for \texttt{RDFW} vs. $\mathcal{O}(\epsilon^{-1})$ for \texttt{RFW}) is the necessary cost of variance control when estimating gradient directions from bit-wise comparisons. Furthermore, while the factor $d$ appears in our complexity, it reflects the "information budget" transition: where \texttt{RFW} requires a $d$-dimensional gradient vector per call, our method requires $d$ individual $1$-bit comparisons, maintaining competitive total information costs.
%     \item \textbf{Comparison with Euclidean Dueling Baselines:} Our \texttt{RDNGD} and \texttt{RRDNGD} algorithms can be viewed as Riemannian generalizations of the NGD methods proposed by \cite{pmlr-v139-saha21b}. While we use a similar gradient estimator, our contribution extends beyond a trivial generalization to the manifold setting. First, we refine the analysis by optimizing the choice of the perturbation parameter $\nu$ and improving the theoretical constants. Second, regarding the problem scope, the original results in \cite{pmlr-v139-saha21b} were strictly limited to the \textbf{Euclidean unconstrained convex} case. In contrast, our framework covers a much broader spectrum: we establish convergence guarantees for \textbf{non-convex} functions in the unconstrained Riemannian setting, and further extend the framework to address \textbf{constrained} g-convex problems on manifolds. Finally, for the strongly g-convex setting using the restart strategy, our modified parameter selection yields a sharper convergence bound: specifically, we improve the complexity's dependence on the initial distance squared $D$ from linear to logarithmic.
% \end{itemize}
% \end{remark}
%\begin{remark}[\textbf{Comparison with Euclidean Dueling Baselines}]\label{rmk:comparison_saha}
%Our \texttt{RDNGD} and \texttt{RRDNGD} algorithms can be viewed as Riemannian generalizations of the NGD methods proposed by \cite{pmlr-v139-saha21b}. While we employ a similar gradient estimator, our contribution extends significantly beyond a trivial generalization to the manifold setting. We highlight four key improvements:

% \begin{itemize}[leftmargin=*]
%     \item \textbf{Refined Theoretical Analysis:} We provide a sharper analysis of the estimator by optimizing the choice of the perturbation parameter $\nu$ (allowing a larger upper bound) and improving theoretical constants. Specifically, we derive a tighter lower bound for the geometric constant $\hat{C}$ and eliminate the logarithmic factor in the probability bound $\gamma_x$, which were present in the Euclidean analysis.
    
%     \item \textbf{Broader Problem Scope:} The original results in \cite{pmlr-v139-saha21b} were strictly limited to the Euclidean unconstrained convex case. In contrast, our framework covers a much broader spectrum, establishing convergence guarantees for non-convex functions in the unconstrained Riemannian setting.
    
%     \item \textbf{Sharper Convergence Bounds:} For the strongly g-convex setting using the restart strategy ({RRDNGD}), our modified parameter selection yields a superior convergence rate. Specifically, we improve the complexity's dependence on the initial squared distance $D$ from linear ($\mathcal{O}(D)$) to logarithmic ($\mathcal{O}(\log D)$).
    
%     \item \textbf{Extension to Projection-Free Optimization:} We further demonstrate the versatility of our estimator by integrating it into the Frank-Wolfe framework. This results in the Riemannian Dueling Frank-Wolfe (RDFW) algorithm, which efficiently solves constrained g-convex problems without requiring computationally expensive projections.
% \end{itemize}
%\end{remark}

% \begin{remark}
% While \cite{RFW} states the constraint requirement as geodesically convex, their analysis relies on the well-definedness of the logarithm map $\text{Exp}_x^{-1}(y)$, which implicitly assumes uniqueness of geodesics within the feasible set. In our work, we explicitly require geodesically uniquely convex sets. This distinction is crucial for general manifolds (especially those with positive curvature), where geodesic convexity does not automatically imply uniqueness. This ensures our direction-finding steps are mathematically rigorous and well-defined.
% \end{remark}





\section{Preliminaries}
We consider smooth Riemannian manifold $\mathcal{M}$ equipped with a Riemannian metric $g$. For any point $x \in \mathcal{M}$, the metric induces an inner product on the tangent space $\mathrm{T}_x \mathcal{M}$, denoted as $\langle \cdot, \cdot \rangle_x$, and an associated norm $\|v\|_x := \sqrt{\langle v, v \rangle_x}$. The Riemannian gradient $\mathrm{grad} f(x) \in \mathrm{T}_x \mathcal{M}$ is uniquely defined as the tangent vector satisfying $\langle \mathrm{grad} f(x), v \rangle_x = \mathrm{D}f(x)[v]$ for all $v \in \mathrm{T}_x \mathcal{M}$, where $\mathrm{D}f(x)[\cdot]$ denotes the differential. To generalize the concept of straight lines in Euclidean space to the manifold, we use geodesics defined as curves that locally minimize distance between two points. The exponential map, $\mathrm{Exp}_x: \mathrm{T}_x \mathcal{M} \to \mathcal{M}$, maps a tangent vector $v$ to $y = \mathrm{Exp}_x(v)$ by following the geodesic starting at $x$ with velocity $v$ for unit time. Conversely, the logarithmic map $\mathrm{Log}_x: \mathcal{M} \to \mathrm{T}_x \mathcal{M}$ is defined as the local inverse of the exponential map, mapping a point $y$ on the manifold back to the tangent vector $v = \mathrm{Log}_x(y)$ such that the geodesic distance between $x$ and $y$ is given by $\|v\|_x$. We use $\Gamma^x_y$ to denote the parallel transport operator that moves a vector along a geodesic while preserving its norm and direction from $\mathrm{T}_y \mathcal{M}$ to $\mathrm{T}_x \mathcal{M}$. 
We refer the reader to Appendix \ref{appendix:prelim} for formal definitions and detailed properties of these concepts.


Now we present some definitions and assumptions. 

\begin{definition}[Geodesic $L$-smoothness]\label{as:lipgrad}
A differentiable function $f : \mathcal{M} \to \mathbb{R}$ is said to be geodesically $L$-smooth if $\rgrad f(x)$ is $L$-Lipschitz, i.e. for any $x, y \in \mathcal{M}$,
\begin{equation*}
\|\rgrad f(x) - \Gamma^x_y \rgrad f(y)\|_x \leq L d(x, y).
\end{equation*}
This also implies, 
\[    |f(y) - f(x) - \langle \rgrad f(x), \Log_x(y) \rangle_x| \leq \frac{L}{2} d^2(x, y).
\]
\end{definition}

\begin{definition}[Geodesically uniquely convex set]\label{def:geodesically_uniquelly_convex_set} 
We call $\mathcal{X}\subseteq \mathcal{M}$ a geodesically uniquely convex set if for any $x,y\in \mathcal{X}$, there exists a unique geodesic $\gamma$ such that 
\begin{align*}
\gamma(0) = x, \quad 
\gamma(1) = y, \quad 
\text{and }~ \gamma(t)\in \mathcal{X}~\text{for all } t\in[0,1].
\end{align*}
\end{definition}

\begin{definition}[Geodesic (strong) convexity]\label{as:strongcon}
Let $\mathcal{X}$ be a geodesically uniquely convex set. The function $f : \mathcal{X}\subseteq\mathcal{M} \to \mathbb{R}$ is geodesically $\alpha$-strongly convex, i.e., for any $x, y \in \mathcal{X}$, it holds that
\begin{equation}\label{geo-convex-inequality}
f(y) \geq f(x) + \langle \rgrad f(x), \Log_x(y) \rangle_x + \frac{\alpha}{2} d^2(x, y).
\end{equation}
If $f$ satisfies \eqref{geo-convex-inequality} with $\alpha=0$, then $f$ is geodesically convex.
\end{definition}
\iffalse
\begin{definition}[Geodesic convexity]\label{as:conv}
Let $\mathcal{X}$ be a geodesically uniquely convex set. The function $f : \mathcal{X}\subseteq\mathcal{M} \to \mathbb{R}$ is geodesically convex if for any $x, y \in \mathcal{X}$, it holds that
\begin{equation}\label{geo-convex-inequality}
f(y) \geq f(x) + \langle \rgrad f(x), \Log_x(y) \rangle_x.
\end{equation}
\end{definition}

\begin{definition}[Geodesic $\alpha$-strong convexity]\label{as:strongcon}
Let $\mathcal{X}$ be a geodesically uniquely convex set. The function $f : \mathcal{X}\subseteq\mathcal{M} \to \mathbb{R}$ is geodesically $\alpha$-strongly convex, i.e., for any $x, y \in \mathcal{X}$, it holds that
\begin{equation*}
f(y) \geq f(x) + \langle \rgrad f(x), \Log_x(y) \rangle_x + \frac{\alpha}{2} d^2(x, y).
\end{equation*}
\end{definition}
\fi
\begin{remark}
We adopt the notion of the geodesically uniquely convex set $\mathcal{X}$ \citep{pmlr-v162-kim22k} to address the ambiguity associated with the possible non-uniqueness of geodesics. 
This relaxes the global property to a local property, enabling a rigorous definition of convexity. This also ensures that $\mathrm{Exp}_x$ is a diffeomorphism for any $x$ in the geodesically uniquely convex set, and thus ensures that $\Log_x(y)=\mathrm{Exp}^{-1}_x(y)$ is well-defined. 


\begin{assumption}\label{as:curvature_lb}
The sectional curvature of $\mathcal{M}$ is bounded below by $\kappa$.
\end{assumption}

\begin{assumption}[Nonexpansive projection]\label{as:proj}
We assume that the set $\mathcal{X}$ is geodesically uniquely convex, and the projection oracle $\mathcal{P}_{\mathcal{X}}:\mathcal{M}\rightarrow \mathcal{X}$ defined as
$
\mathcal{P}_{\mathcal{X}}(x) := \{ y \in \mathcal{X} : d(x, y) = \inf_{z \in \mathcal{X}} d(x, z) \},
$ is nonexpansive, i.e., $d(\mathcal{P}_{\mathcal{X}}(x), \mathcal{P}_{\mathcal{X}}(y)) \leq d(x, y)$ for all $x,y\in\mathcal{M}$. 
\end{assumption}
\end{remark}



\section{Riemannian Dueling Normalized Gradient Descent Method}\label{sec:section3}
In this section, we present our Riemannian Dueling Normalized Gradient Descent (RDNGD) method for solving \eqref{eq:target1.1}. For ease of notation, we use $\langle\cdot,\cdot\rangle$ to denote the inner product $\langle\cdot,\cdot\rangle_x$ on $\mathrm{T}_x\mathcal{M}$ and $\|\cdot\|$ to denote $\|\cdot\|_x$, when there is no ambiguity. For a tangent space $\mathrm{T}_x\mathcal{M}$ with an orthonormal basis $\{e_1,e_2,\ldots e_d\}$, any $v\in \mathrm{T}_x\mathcal{M}$ can be expressed as $v=\sum_{i=1}^d v_i e_i$ for some $v_1,v_2,\ldots,v_d\in\mathbb{R}$. When the choice of basis is clear from context, we represent the tangent vector by its coefficient vector $ v := [v_1, v_2, \ldots, v_d]^\top$.
For a given point $x\in\mathcal{M}$, we use $\mathcal{S}_{\mathrm{T}_x\mathcal{M}}(r)$ to denote the sphere with center $x$ and radius $r$ on $\mathrm{T}_x\mathcal{M}$. For a set $S$, let $\mathrm{Unif}(S)$ represent the uniform distribution on $S$. Sampling a tangent vector uniformly from $\mathcal{S}_{\mathrm{T}_x\mathcal{M}}(1)$ is equivalent with sampling the coefficient vector $u$ uniformly from $\mathcal{S}_d(1):=\{x\in \mathbb{R}^d:\|x\|_2=1\}$. 

We now introduce our Riemannian gradient direction estimator using dueling oracles $\mathcal{Q}_f(x,y)$. Although such estimators have been proposed for the Euclidean space \cite{zhang2024comparisons,pmlr-v139-saha21b}, none of them applies to the Riemannian setting directly. 
Our Riemannian gradient direction estimator is defined as:
\begin{align}\label{eq:graddirestimator}
h_\pert(x)=\mathcal{Q}_f(\mathrm{Exp}_x(\pert u), \mathrm{Exp}_x( -\pert u))u,
\end{align}
where $\pert>0$ is the perturbation radius and $u\sim \mathrm{Unif}(\mathcal{S}_{\mathrm{T}_x\mathcal{M}}(1))$. The exponential map ensures that the perturbed points stay on $\mathcal{M}$. 
$u\sim \mathrm{Unif}(\mathcal{S}_{\mathrm{T}_x\mathcal{M}}(1))$ is sampled by projecting a random Gaussian vector onto the tangent space  followed by normalization. 
The following results establish the relation between $h_\nu(x)$ in \eqref{eq:graddirestimator} and the normalized gradient. In Lemma~\ref{lem:grad_with_bias} we first show that $h_\pert(x)$ coincides with $\operatorname{sign}(\langle \rgrad f(x),u\rangle_x)u$  with high probability.

\begin{lemma}\label{lem:grad_with_bias}
Assume $f$ is geodesically $L$-smooth. Let $u \sim \mathrm{Unif(\mathcal{S}_{\mathrm{T}_x\mathcal{M}}(1))}$, and $\pert \in (0,1)$. Then with probability at least $1 - \gamma_x$, we have 
\begin{equation*}
h_{\pert}(x)
= \operatorname{sign}(\langle\rgrad f(x), u\rangle_x)  u,
\end{equation*}
where $h_{\pert}(x)$ is defined in \eqref{eq:graddirestimator} and 
\begin{equation}\label{def-gamma_x}
\gamma_x = \sqrt{\frac{d}{2\pi}}\frac{{L}\pert}{\norm{\rgrad f(x)}}
\end{equation}
\end{lemma}
Now in Lemma~\ref{lem:estimating_normalized_gradient}, setting $v=\rgrad f(x)$, we show that $ \frac{\sqrt{d}}{\hat{C}}\operatorname{sign}(\langle \rgrad f(x),u\rangle_x)u$ is an unbiased estimator of the normalized gradient for some universal constant $\hat{C}$. 
\begin{lemma}\label{lem:estimating_normalized_gradient}
Fix $x\in\mathcal{M}$ and a nonzero vector $v\in \mathrm{T}_x\mathcal{M}$. 
Let $u\sim\mathrm{Unif}(\mathcal{S}_{\mathrm{T}_x\mathcal{M}}(1))$.  
The following holds for some universal constant $\hat{C} \in \big[\frac{1}{\sqrt{2\pi}},1\big]$:
\begin{align}\label{lem:estimating_normalized_gradient-eq}
    \mathbb{E}_u\big[\operatorname{sign}(\langle v,u\rangle_x)u\big]
= \frac{\hat{C}}{\sqrt{d}}\frac{v}{\|v\|_x}.
\end{align}
\end{lemma}
Using Lemma~\ref{lem:grad_with_bias}, and Lemma~\ref{lem:estimating_normalized_gradient}, Proposition \ref{prop:estimate_biased_NG} below shows that the estimator $h_\pert(x)$ is, on average, approximately aligned with the gradient direction.
\begin{proposition}\label{prop:estimate_biased_NG}
Assume $f$ is geodesically $L$-smooth. For any $x\in\mathcal{M}$, $\nu\in (0,1)$ and $v\in S_{\mathrm{T}_x\mathcal{M}}(1)$, it holds that 
\begin{equation*}
\abs{\mathbb{E}_u\bigl[ \langle h_\pert(x), v\rangle\bigr]-\frac{\hat{C}}{\sqrt{d}}
\left\langle\frac{\rgrad f(x)}{\|\rgrad f(x)\|},v\right\rangle_x }
\le 2\gamma_x,
\end{equation*}
for some universal constant $\hat{C}\in[\frac{1}{\sqrt{2\pi}},1]$, where $\gamma_x$ is defined in \eqref{def-gamma_x}. 
\end{proposition}
\begin{remark}
Setting $v=\rgrad f(x)/\|\rgrad f(x)\|$, and $\pert=\sqrt{2\pi}/d$, we see that $\sqrt{d}h_\pert (x)/\hat{C}$ is approximately aligned with gradient direction with error $O(L/(\hat{C}\|\rgrad f(x)\|)$. This bias is unavoidable as $\hat{C}$, and $\|\rgrad f(x)\|$ are unknown. This is a key distinction of our setting from standard stochastic Riemannian optimization algorithms where the latter relies on unbiased gradient estimators \cite{SGD_riemannian, tripuraneni2018averaging}. Lemma~\ref{lem:estimating_normalized_gradient} improves the lower bound on $\hat{C}$ from $\tfrac{1}{20}$ in \citep{pmlr-v139-saha21b} to $\tfrac{1}{\sqrt{2\pi}}\approx0.4$, leading to up to an eightfold reduction in gradient direction estimation error bound.
Lemma~\ref{lem:grad_with_bias} sharpens the bias bound by removing the logarithmic factor $\sqrt{\log\big(|\nabla f(\mathbf{x})|/(\sqrt{d}L\nu)\big)}$ present in $\gamma_x$ in \citep{pmlr-v139-saha21b} for the Euclidean case.
\end{remark}

Building on the estimator $h_\pert(x)$, we now present our RDNGD method (Algorithm \ref{alg:alg1}) for solving \eqref{eq:target1.1}. Inputs of RDNGD are the initial point $x_0\in\mathcal{X}\subseteq\mathcal{M}$, learning rate $\{\eta_k\}$, perturbation radius $\nu$, and maximum iteration number $T>1$. $\hat{x}_k$ records the best iterate in the first $k$ iterations. 
The RDNGD method applies to two cases: (i) unconstrained problem where $f$ is geodesically $L$-smooth and $\mathcal{X}=\mathcal{M}$; (ii) constrained convex problem where $f$ is geodesically convex and geodesically $L$-smooth. The iteration complexities of RDNGD for obtaining an $\epsilon$-stationary solution in cases (i) and (ii) are presented in Theorem \ref{thm:beta_smoothness_complexityapp}, and Theorem \ref{thm:beta_smoothness_convex_complexityapp} respectively. For both cases, we consider two different choices of step sizes: the constant step size and the cosine annealing step size. The latter is defined as follows. 

\begin{algorithm}[t]
\caption{Riemannian Dueling Normalized Gradient Descent (\textbf{RDNGD})}
\label{alg:alg1}
\begin{algorithmic}[1]
\STATE Set $\hat{x}_0=x_0$
\FOR{$k = 0, 1, \dots, T-1$}
    \STATE Sample $u_k \sim \mathrm{Unif}(\mathcal{S}_{\mathrm{T}_{x_k}\mathcal{M}}(1))$
    \STATE $h_\pert(x_k) = \mathcal{Q}_f(\mathrm{Exp}_{x_k}(\pert u_k),\mathrm{Exp}_{x_k}(-\pert u_k))u_k$
    \STATE Update $x_{k+1} = \mathcal{P}_{\mathcal{X}}\left(\mathrm{Exp}_{x_k}(- \eta_k h_\pert(x_k))\right)$
    \STATE $\hat{x}_{k+1}=b_k\hat{x}_{k}+(1-b_k)x_{k+1}$, 
    \STATE \quad where $b_k=(\mathcal{Q}(x_{k+1},\hat{x}_k)+1)/2$
    %\mathbb{1}(f(x_{k+1})>f(\hat{x}_k))
\ENDFOR
\STATE \textbf{Return} $\hat{x}:=\hat{x}_{T}$
\end{algorithmic}
\end{algorithm}

\begin{definition}[\bf Cosine Annealing step size]\label{def:cosine_annealing_step size}
Given the initial step size $\eta_0$, the minimum step size $\eta_{\min}\in [0,\eta_0]$, and the total number of iterations $T>0$, the cosine annealing step size at the $k$-th iteration is defined as:
\begin{equation}\label{eq:cosine-anneal}
\eta_k = \eta_{\min} + \frac{1}{2}(\eta_{0}-\eta_{\min})\left(1+\cos\left(\frac{k\pi}{T}\right)\right)
\end{equation}
for $k=0,1,\ldots,T$.
\end{definition}

\begin{theorem}\label{thm:beta_smoothness_complexityapp}
Assume $f$ is geodesically $L$-smooth. We consider using RDNGD (Algorithm~\ref{alg:alg1}) to solve \eqref{eq:target1.1} with $\mathcal{X}=\mathcal{M}$\footnote{Note that in this case, the projection operator in Algorithm \ref{alg:alg1} can be ignored.}. For any given $\epsilon>0$, RDNGD returns an $\epsilon$-stationary point with the following two sets of inputs. 
\vspace{-0.05in}
\begin{itemize}[leftmargin=*]
\item[(i)] \textbf{Constant step size.}
\[    T = \left\lceil \frac{L^2d(D+1)^2}{\hat{C}^2\epsilon^2} \right\rceil,\
    \eta_k\equiv \eta=\frac{1}{\sqrt{T}}, \
    \pert = \frac{\hat{C}\sqrt{2\pi}}{4dL}\epsilon.
\]
In this case, RDNGD returns a point $x_R$ such that $\mathbb{E}[\|\rgrad f(x_R)\|] < \epsilon$
where $x_R$ is chosen uniformly at random from $\{x_0, \dots, x_{T-1}\}$. %with probability $\mathbb{P}(R=r)=\frac{1}{T}$. 
 Throughout this paper, $D:= d^2(x_0,x^*)$ denotes the squared distance between the initial point and the optimal point. 

\item[(ii)] \textbf{Cosine annealing step size.}
\begin{equation}\label{thm:beta_smoothness_complexityapp-cosine-parameter}
\begin{split}
    T &= \left\lceil \left( \frac{2\sqrt{d}}{\epsilon \hat{C}} \left( 2L D + \tfrac{1}{2}L \right) \right)^{2} \right\rceil,\quad \eta_0 = \frac{1}{\sqrt{T}},\\
    &\eta_k \text{ defined in } \eqref{eq:cosine-anneal}, \quad
    \pert = \frac{\hat{C}\sqrt{\pi}}{2\sqrt{2}dL}\epsilon.
\end{split}
\end{equation}
In this case, RDNGD returns a point $x_R$ such that $\mathbb{E}[\|\rgrad f(x_R)\|] < \epsilon$
% \begin{equation}\label{thm:beta_smoothness_complexityapp-cosine-parameter-conclusion}
% \mathbb{E}[\|\rgrad f(x_R)\|] < \epsilon,
% \end{equation}
where $x_R$ is chosen at random from $\{x_0, \dots, x_{T-1}\}$ with probability $\mathbb{P}(R=r)=\frac{\eta_r}{\sum_{k=1}^T\eta_k}$.
\end{itemize}
\end{theorem}

\begin{theorem}\label{thm:beta_smoothness_convex_complexityapp}
Assume $f$ is geodesically $L$-smooth and geodesically convex, and Assumptions \ref{as:curvature_lb} and \ref{as:proj} hold. Consider using RDNGD (Algorithm \ref{alg:alg1}) to solve \eqref{eq:target1.1} over a bounded set $\mathcal{X}\subset\mathcal{M}$ of diameter $\mathfrak{D}$. For any $\epsilon>0$, RDNGD returns an $\epsilon$-optimal solution satisfying $\mathbb{E}[f(\hat{x}_{T})] - f(x^*) < \epsilon$ with the following two sets of inputs.
\begin{itemize}[leftmargin=*]
\item[(i)] \textbf{Constant step size.}
\begin{equation}\label{thm:beta_smoothness_convex_complexityapp-constant-step-param}
\begin{split}
    T &= \left\lceil1+\frac{16\pi d{L}\bar{\zeta}}{\epsilon}D\right\rceil, \quad 
    \eta =\frac{\sqrt{\epsilon}}{4\sqrt{\pi}\bar{\zeta}\sqrt{d{L}}}, \\
    \pert &= \frac{1}{4\sqrt{2} d}\frac{(\epsilon/{L})^{3/2}}{(\sqrt{D}+\eta T)^2}.
\end{split}
\end{equation}
Here and in the rest of the paper, $\bar{\zeta} :=  {\sqrt{|\kappa|}\mathfrak{D}}/{\tanh(\sqrt{|\kappa|} \mathfrak{D})}$, where $\kappa$ is the sectional curvature lower bound in Assumption \ref{as:curvature_lb}.

\item[(ii)] \textbf{Cosine annealing step size.}
\begin{equation}\label{thm:beta_smoothness_convex_complexityapp-cosine-step-param}
\begin{split}
    T &= \left\lceil 1 + \frac{32\pi d{L}\bar{\zeta}}{\epsilon}D \right\rceil, \quad 
    \eta_0 = \frac{\sqrt{\epsilon}}{4\sqrt{\pi}\bar{\zeta}\sqrt{d{L}}}, \\
    \pert &= \frac{1}{4\sqrt{2} d}\frac{(\epsilon/{L})^{3/2}}{(\sqrt{D}+\frac{\eta_0}{2}T)^2}.
\end{split}
\end{equation}
\end{itemize}
\end{theorem}


\begin{remark}
Here we remark that the Assumptions \ref{as:curvature_lb} and \ref{as:proj} naturally hold if $\mathcal{M}$ is a Hadamard manifold, and they are standard assumptions under such scenario. See \citep{Bacak2014Hadamard,zhang2016firstordermethodsgeodesicallyconvex, li2024zeroth}.
\end{remark}

\begin{remark}
Our rates in
Theorem~\ref{thm:beta_smoothness_convex_complexityapp} match the rates in \citep{pmlr-v139-saha21b} when the problem is reduced to the Euclidean case. However, \citep{pmlr-v139-saha21b} requires strictly smaller choices of step size to ensure convergence due to logarithmic factors in the analysis, whereas our analysis permits larger step size potentially leading to faster convergence in practice. Furthermore, when reduced to Euclidean case, our results in Theorems \ref{thm:beta_smoothness_complexityapp} and \ref{thm:beta_smoothness_convex_complexityapp} also improve the dimension dependence in \citep{zhang2024comparisons} by a factor of $\log d$. Additionally, our geometry-aware bounds are adaptive to the intrinsic manifold structure, in contrast to the ambient-dimension–dependent factors in prior analyses.
\end{remark}


We now introduce a new algorithm that achieves linear convergence rate for solving \eqref{eq:target1.1} when $f$ is geodesically $\alpha$-strongly convex and geodesically $L$-smooth. The key observation is that under strong convexity, controlling $f(x)-f(x^*)$ directly controls the squared distance $d^2(x,x^*)$. We exploit this idea by designing an algorithm that runs in phases: in phase $k$, we run our RDNGD (Algorithm \ref{alg:alg1}) to obtain an iterate with pre-given function value sub-optimality $\epsilon_k$, which leads to reductions in the estimation error. By reducing $\epsilon_k$ in each iteration, the convergence rate is improved to linear rate. 
The algorithm is called Riemannian Recurrent Dueling Normalized Gradient Descent (RRDNGD) and is presented in Algorithm~\ref{alg:alg2}, where we used a constant phase length $t = \lceil{1+{64\pi d L \bar{\zeta}}/{\alpha}}\rceil$. 


\begin{algorithm}[t]
\caption{Riemannian Recurrent Dueling Normalized Gradient Descent (\textbf{RRDNGD})}
\label{alg:alg2}
\begin{algorithmic}[1]
\STATE \textbf{Input:} Initial point $x_0 \in \mathcal{X} \subset \mathcal{M}$, constant phase length $t = \left\lceil 1 + \frac{64\pi d L \bar{\zeta}}{\alpha} \right\rceil$, total number of phases $K$
\STATE \textbf{Initialize:} $x^{(0)}=x_0$, $D_0 = D$
\FOR{$k = 0, 1, \dots, K-1$}
    \STATE Run RDNGD with inputs $(x^{(k)}, \eta_k, \nu_k, t)$ to obtain $x^{(k+1)}$
    \STATE Update $D_{k+1} = D_{k}/2$
\ENDFOR
\STATE \textbf{Return} $\hat{x} = x^{(K)}$
\end{algorithmic}
\end{algorithm}


We now present the oracle complexity of RRDNGD. 
\begin{theorem}\label{thm:strong_convexity_complexity}
Assume $f$ is geodesically $L$-smooth and geodesically $\alpha$-strongly convex. Assume Assumptions \ref{as:curvature_lb} and \ref{as:proj} hold, and the global optima $x^*$ lies in the interior of $\mathcal{X}$. Consider using RRDNGD (Algorithm \ref{alg:alg2}) to solve \eqref{eq:target1.1} over a bounded set $\mathcal{X}\subset\mathcal{M}$ of diameter $\mathfrak{D}$. For any $\epsilon>0$, we set $K = \left\lceil \log_2\left( \frac{l(L,\mathfrak{D})^2 D}{\epsilon^2} \right) \right\rceil$, where $l(L,\mathfrak{D})$ is constant dependent on $L$, and $\mathfrak{D}$. In the $k$-th iteration, set:
\[    \epsilon_k = \frac{\alpha D}{2^{2-k}}, 
    \eta_k = \frac{1}{4\bar{\zeta}}\sqrt{\frac{\epsilon_k}{\pi dL}}, 
    \nu_k = \frac{1}{4\sqrt{2} d}\frac{(\epsilon_k/L)^{3/2}}{(2\sqrt{\epsilon_{k}/\alpha}+\eta_k t)^2}.
\]
Then RRDNGD returns an $\epsilon$-optimal solution $\hat{x}$ satisfying $\mathbb{E}[f(\hat{x})] - f(x^*) \leq \epsilon$
with an oracle complexity of
$\mathcal{O}\left( \frac{dL\bar{\zeta}}{\alpha} \log\left( \frac{l(L,\mathfrak{D}) \sqrt{D}}{\epsilon} \right) \right)$\footnote{Our complexity improves the dependency on the initial squared distance from linear in \citep{pmlr-v139-saha21b}, which is $\mathcal{O}\big( \frac{dL}{\alpha} ( \log \frac{\alpha}{\epsilon} + \|\mathbf{x}_1 - \mathbf{x}^*\|^2 ) \big)$, to logarithmic.}.
\end{theorem}

\iffalse
\begin{theorem}\label{thm:strong_convexity_complexity}
Assume $f$ is geodesically $L$-smooth and geodesically $\alpha$-strongly convex, and Assumptions \ref{as:curvature_lb} and \ref{as:proj} hold. Consider using RRDNGD (Algorithm \ref{alg:alg2}) to solve \eqref{eq:target1.1} over a bounded set $\mathcal{X}\subset\mathcal{M}$ of diameter $\mathfrak{D}$. For any $\epsilon>0$, we set $K = \left\lceil \log_2\left( \frac{l^2 D}{\epsilon^2} \right) \right\rceil$, where $l$ denotes the Lipschitz constant of $f$ on $\mathcal{X}$. We set parameters used in the $k$-th iteration as:{\color{red}define $D_k$}
\begin{align}
    \epsilon_k = \frac{\alpha D_k}{4},  
    \eta_k = \frac{1}{4\bar{\zeta}}\sqrt{\frac{\epsilon_k}{\pi dL}},  
    \nu_k = \frac{1}{4\sqrt{2} d}\frac{(\epsilon_k/L)^{3/2}}{(\sqrt{D_{k}}+\eta_k t)^2}.
\end{align}
Then RRDNGD returns an $\epsilon$-optimal solution $\hat{x}$ satisfying $\mathbb{E}[f(\hat{x})] - f(x^*) \leq \epsilon$
with an oracle complexity of
$\mathcal{O}\left( \frac{dL\bar{\zeta}}{\alpha} \log\left( \frac{l \sqrt{D}}{\epsilon} \right) \right)$\footnote{Our complexity improves the dependency on the initial squared distance from linear in \citep{pmlr-v139-saha21b}, which is $\mathcal{O}\big( \frac{dL}{\alpha} ( \log \frac{\alpha}{\epsilon} + \|\mathbf{x}_1 - \mathbf{x}^*\|^2 ) \big)$, to logarithmic.}.
\end{theorem}

\begin{remark}
When $f$ is geodesically $L$-smooth and the domain $\mathcal{X}$ is bounded, $f$ is geodesically Lipschitz continuous on $\mathcal{X}$ with some Lipschitz constant $l$.
\end{remark}
\fi
\section{Riemannian Dueling Frank-Wolfe Method}
Note that RDNGD method requires a projection in each iteration, which can be expensive for some applications. To overcome this issue, in this section, we design a projection-free Riemannian dueling Frank-Wolfe (RDFW) method using comparison oracles. In the Euclidean space, the Frank-Wolfe method calls a linear minimization oracle in each iteration. On manifolds, at iteration $k$, Frank-Wolfe method requires minimizing the following analogous subproblem:
\begin{equation}\label{eq:rfwsubprob}
 \arg\min_{z\in\mathcal{X}} \left\langle \rgrad f(x_k), \Log_{x_k}(z)\right\rangle.
\end{equation}
In some applications, unlike projection, \eqref{eq:rfwsubprob} admits a closed–form solution. % that are inexpensive to evaluate. Motivated by this property, we develop a Riemannian Frank--Wolfe method driven by a dueling oracle.
For example, consider the manifold of symmetric positive definite (SPD) matrices $\mathcal{M}=\mathcal{S}_{++}^n:=\{X\in\mathbb{R}^{n\times n}\mid X^T=X, X\succ 0\}$ with a geodesically convex ``interval'' constraint
$\mathcal{X} := \{ X \in \mathcal{M} : L \preceq X \preceq U \}$ for SPD matrices $L$ and $U$. When $L$ and $U$ do not commute, the Riemannian projection 
\iffalse
\begin{align*}
X_{\mathrm{proj}}
:=
\arg\min_{Y\in\mathcal{X}}
 d_{\mathcal{M}}(Y,X)
\end{align*}
\fi
requires an expensive iterative solver. In contrast, \eqref{eq:rfwsubprob} admits a closed–form solution that is cheap to compute \citep{RFW}.
%\subsection{Why variance control is necessary for FW}
\begin{algorithm}[t]
\caption{Riemannian Dueling Frank-Wolfe Method (\textbf{RDFW})}
\label{alg:RFW_batch}
\begin{algorithmic}[1]
\STATE \textbf{Input:} Initial point $x_0 \in \mathcal{X} \subseteq \mathcal{M}$, total number of iterations $T$
\FOR{$k=0, 1, \dots, T-1$}
    \FOR{$j=1, \dots, M_k$}
        \STATE Sample $u^{(k)}_j \sim \mathrm{Unif}(\mathcal{S}_{\mathrm{T}_{x_k}\mathcal{M}}(1))$
        \STATE $o^{(k)}_j=\mathcal{Q}_f(\mathrm{Exp}_{x_k}(\pert_k u^{(k)}_j),\mathrm{Exp}_{x_k}(-\pert_k u^{(k)}_j))$
        \STATE $h_{\pert_k}^j(x_k) = o_j^{(k)} u^{(k)}_j$
    \ENDFOR
    \STATE $\bar h_k \gets \frac{1}{M_k}\sum_{j=1}^{M_k} h_{\pert_k}^j(x_k)$
    \STATE $z_k \gets \arg\min_{z\in\mathcal{X}} \left\langle \bar h_k, \Log_{x_k}(z)\right\rangle$
    \STATE $s_k \gets \frac{2}{k+3}$
    \STATE $x_{k+1}=\mathrm{Exp}_{x_k}(s_k\Log_{x_k}(z_k))$
\ENDFOR
\STATE \textbf{Return} $\hat{x} = x_{T}$
\end{algorithmic}
\end{algorithm}

We now present our RDFW algorithm with the detailed description in Algorithm \ref{alg:RFW_batch}. In the absence of gradient information in dueling feedback setting, RDFW solves the subproblem~\eqref{eq:rfwsubprob} with $\rgrad f(x)$ replaced by our gradient direction estimator. However, the solution to \eqref{eq:rfwsubprob} is highly sensitive to the noise in gradient estimation (see Section~\ref{sec:fw_sensitivity_example}). To reduce the noise variance, in the $k$-th iteration, we use the batch gradient direction estimator given by:
\begin{equation}\label{eq:batchgrad}
\bar{h}_k = \frac{1}{M_k}\sum_{j=1}^{M_k} o^{(k)}_ju^{(k)}_j,
\end{equation}
where we use $M_k$ i.i.d. directions $u^{(k)}_j \sim \mathrm{Unif}(\mathcal{S}_{\mathrm{T}_{x_k}\mathcal{M}}(1))$. The variance of $\bar{h}_k$ reduces at a $\frac{1}{M_k}$ rate. We now solve the following Riemannian Frank–Wolfe subproblem:
\begin{align*}
z_k \in \arg\min_{z\in\mathcal{X}} \left\langle \bar{h}_k, \Log_{x_k}(z)\right\rangle.
\end{align*}
%This step avoids the expensive projections, and 
The update $x_{k+1}$ is obtained by moving along the geodesic from $x_k$ to $z_k$, and this step preserves feasibility. Theorem \ref{thm:RFW_batch} gives the oracle complexity of RDFW.  
\begin{theorem}\label{thm:RFW_batch}
Assume $f$ is geodesically $L$-smooth and geodesically convex. For any given $\epsilon>0$, suppose RDFW (Algorithm~\ref{alg:RFW_batch}) is run with 
\begin{align*}
    &\pert_k=\sqrt{\frac{\pi}{8}}\frac{\hat{C}}{dL\mathfrak{D}}\frac{2}{k+3}, \quad M_k=\frac{8dL\mathfrak{D}^2}{\hat{C}^2}(k+3),\\
    &T=\left\lceil \frac{2(f(x_{0})-f(x^*)) + 4L\mathfrak{D}^2+12}{\epsilon}\right\rceil,
\end{align*} 
then $\mathbb{E}\left[f(x_{T}) - f(x^*)\right]<\epsilon$, and the oracle complexity is $2\sum_{k=0}^{T-1} M_k =\mathcal{O}(d/\epsilon^{2})$.
\end{theorem}

\iffalse
\begin{remark}
A critical distinction between our setting and standard Stochastic Frank-Wolfe (SFW) algorithms lies in the nature of the gradient estimator. Standard SFW typically assumes access to unbiased gradient estimators. In contrast, our proposed zeroth-order estimator is inherently biased. While this bias is controlled by the smoothing radius, a single-sample estimator suffers from high variance. Consequently, we employ a batching strategy to explicitly reduce this variance, ensuring that variance is sufficiently bounded to guarantee convergence.
\end{remark}
\fi
\section{Numerical Experiments}

In this section, we conduct numerical experiments to test our algorithms on both synthetic data and real applications.

\subsection{Synthetic Problems}
We consider three problems with synthetic data: Rayleigh quotient maximization, the Karcher mean problem, and the Karcher mean problem with an additional constraint. 

\subsubsection{Rayleigh Quotient Maximization}\label{sec:rayleigh}

Given a symmetric matrix $A \in \mathbb{R}^{d \times d}$, the Rayleigh quotient maximization problem is formulated as the minimization of the negative quadratic form over the unit sphere:
\begin{equation}\label{prob:rayleigh}
    \min_{x \in \mathcal{S}_d(1)} f(x) := -\frac{1}{2} x^\top A x.
\end{equation}
The objective $f$ is geodesically $L$-smooth with $L = \lambda_{\max}(A) - \lambda_{\min}(A)$ \citep{pmlr-v162-kim22k} and $f^* = -\frac{1}{2}\lambda_{\max}(A)$. Following the setup in \cite{pmlr-v162-kim22k}, we generate a random matrix $B \in \mathbb{R}^{d \times d}$ with entries sampled i.i.d.\ from $\mathcal{N}(0, 1/d)$ and set $A = \frac{1}{2}(B + B^\top)$. 
We compare our RDNGD (Algorithm \ref{alg:alg1}) 
with the zeroth-order Riemannian gradient descent method (ZO-RGD) \citep{li2023stochastic}. While ZO-RGD requires function values, RDNGD relies only on comparison oracles yet achieves comparable performance. The algorithms are initialized at $x_0 \sim \mathrm{Unif}(\mathcal{S}_d(1))$. RDNGD-constant uses a constant step size $\eta = 10^{-2}$, and RDNGD-cosine employs cosine annealed step size with $\eta_{0}=10^{-1}$ and $\eta_{\min}=10^{-8}$. 
Both RDNGD variants use $\nu=10^{-8}$. For ZO-RGD \citep{li2023stochastic}, we set the smoothing parameter as $10^{-6}$ and the step size as $\frac{1}{2dL}$. These values are tuned via a grid search. All algorithms are run for $T=50,000$ iterations. 

Figure~\ref{fig:rayleigh} reports $f(x_k)-f(x^*)$ for different methods for \eqref{prob:rayleigh} for dimensions $d=100$ and $d=150$. We see that RDNGD-cosine achieves similar performance as ZO-RGD although the former uses only the comparison oracle. Although RDNGD-constant performed worse than RDNGD-cosine and ZO-RGD but still achieves acceptable accuracy.
\begin{figure}[t]
    \centering
    \begin{minipage}[b]{0.48\linewidth}
        \centering
        \includegraphics[width=\linewidth]{results/rayleigh_final_d=100.png}
        \centerline{\small (a) d=100}
    \end{minipage}
    \hfill
    \begin{minipage}[b]{0.48\linewidth}
        \centering
        \includegraphics[width=\linewidth]{results/rayleigh_final_d=150.png}
        \centerline{\small (b) d=150}
    \end{minipage}
    \caption{Numerical results for Rayleigh quotient maximization.}
    \label{fig:rayleigh}
\end{figure}
\vspace{-0.1in}
\subsubsection{Karcher Mean Problem}
In this problem, the goal is to find the geometric mean or Karcher mean of $m$ SPD matrices $\{A_i\}_{i=1}^m \subset \mathcal{S}_{++}^n$, and can be cast as follows:
\begin{equation}\label{prob:Karcher}
    \min_{X \in \mathcal{S}_{++}^n} f(X) := \frac{1}{2m} \sum_{i=1}^m d^2(X, A_i),
\end{equation}
where $d(\cdot,\cdot)$ is the Riemannian distance induced by the affine-invariant metric on $\mathcal{S}_{++}^n$. $f$ in \eqref{prob:Karcher} is geodesically $L$-smooth and geodesically convex on $\mathcal{S}_{++}^n$ \cite{zhang2016firstordermethodsgeodesicallyconvex}. Here we again compare our RDNGD algorithm with the ZO-RGD algorithm. The matrices $A_i$ were randomly generated using the Matrix Mean Toolbox \citep{BINI20131700} with a condition number of $10^6$ and sample size $m=50$. We consider dimensions $n \in \{5, 10\}$. The algorithms are initialized at $X_0$, which is set to the arithmetic mean of all $A_i$'s.
To evaluate the performance, the optimal solution $X^*$ of \eqref{prob:Karcher} is computed using the Richardson-like linear gradient descent algorithm provided in the Matrix Mean Toolbox. RDNGD-constant uses a constant step size $\eta = 10^{-2}$, and RDNGD-cosine employs $\eta_{0}=10^{-1}$ and $\eta_{\min}=10^{-6}$. 
Both RDNGD variants use $\nu=10^{-8}$. For ZO-RGD \citep{li2023stochastic}, we set the smoothing parameter as $10^{-6}$ and the step size as $\frac{1}{20nL}$ and these values are again tuned via a grid search to produce the best results for this problem. All algorithms are run for $T=5,000$ iterations.

\begin{figure}[t]
    \centering
    \begin{minipage}[b]{0.48\linewidth}
        \centering
        \includegraphics[width=\linewidth]{results/spd_final_d=5.png}
        \centerline{\small (a) n=5}
    \end{minipage}
    \hfill
    \begin{minipage}[b]{0.48\linewidth}
        \centering
        \includegraphics[width=\linewidth]{results/spd_final_d=10.png}
        \centerline{\small (b) n=10}
    \end{minipage}
    \caption{Numerical result for Karcher mean problem.}
    \label{fig:Karcher_mean_unconstrained}
\end{figure}
Figure~\ref{fig:Karcher_mean_unconstrained} reports the convergence of different methods on \eqref{prob:Karcher} for dimensions $n\in\{5,10\}$. We see that RDNGD-cosine can achieve comparable accuracy with ZO-RGD. RDNGD-constant performed slightly worse but the accuracy is still on an acceptable level. This is again remarkable because RDNGD can only access a comparison oracle. 

\subsubsection{Karcher Mean with Constraints}
\begin{figure}[t]
    \centering
    \begin{minipage}[b]{0.42\linewidth}
        \centering
        \includegraphics[width=\linewidth]{results/Karcher_Mean_constrained_d=5.png}
        \centerline{\small (a) n=5}
    \end{minipage}
    \hfill
    \begin{minipage}[b]{0.42\linewidth}
        \centering
        \includegraphics[width=\linewidth]{results/Karcher_Mean_constrained_d=10.png}
        \centerline{\small (b) n=10}
    \end{minipage}
    \caption{Numerical result for constrained Karcher mean problem.}
    \label{fig:Karcher_mean_constrained}
\end{figure}
In this subsection, we apply our RDFW method (Algorithm \ref{alg:RFW_batch}) to solve the Karcher mean problem with the constraint $H \preceq X \preceq A$ where $H$ is the harmonic mean, and $A$ is the arithmetic mean \citep{RFW}.
It is a natural constraint as the Karcher mean $G$ satisfies $H \preceq G \preceq A$ \citep{bhatia2007}. Formally, we obtain the following formulation:
\begin{align*}
    \min_{X \in \mathcal{S}_{++}^n,\ H \preceq X \preceq A} ~~~ f(X) := \frac{1}{2m} \sum_{i=1}^m d^2(X, A_i).
\end{align*}
We set $m=50$ and consider dimensions $n \in \{5, 10\}$. The algorithm is initialized at $X_0 = A$. We set the perturbation parameter to $\pert_k=10^{-8}$, the batch size to $M_k=\lceil (k+1)/500\rceil$, and the total number of iterations to $T=5,000$.

Figure~\ref{fig:Karcher_mean_constrained} shows that RDFW converges to high-accuracy solutions for the constrained Karcher mean problem. To our knowledge, RDFW is the first projection-free algorithm to solve this problem using only comparison oracles.
\subsection{Real Applications}
\subsubsection{Attack on Deep Neural Network}
We evaluate our dueling Riemannian attack on the black-box benchmark~\citep{li2023stochastic}, attacking a VGG network on CIFAR-10 under an $\ell_2$-norm constraint. The perturbation is confined to a sphere (dimension $\approx 1000$) centered at the original image $x \in \mathbb{R}^d$. We adopt the exact experimental settings from ~\citep{li2023stochastic}.\footnote{All attacks are implemented by extending the open-source codebase available at \url{https://github.com/JasonJiaxiangLi/Zeroth-order-Riemannian}}For our RDNGD algorithm, we set $\nu=10^{-6}$, step size $\eta=10^{-6}$, and $T=1,000$. We use the batch gradient direction estimator $\bar{h}_k$ \eqref{eq:batchgrad} with batch size of $10$ to mitigate the high variance of the gradient direction estimator in such a high-dimensional space. In contrast, ZO-RGD request 500 samples.
\begin{figure}[t]
    \centering
    % --------------------------------------------------
    % --------------------------------------------------
    \begin{subfigure}[b]{1.0\linewidth}
        \centering
        \setlength{\tabcolsep}{1pt}
        \renewcommand{\arraystretch}{1.0}
        \begin{tabular}{ccccc}
            \includegraphics[width=0.19\linewidth]{results/CIFAR_result/1_Original_cat.png} &
            \includegraphics[width=0.19\linewidth]{results/CIFAR_result/1_PGD_ship.png} &
            \includegraphics[width=0.19\linewidth]{results/CIFAR_result/1_RGD_ship.png} &
            \includegraphics[width=0.19\linewidth]{results/CIFAR_result/1_RZO_ship.png} &
            \includegraphics[width=0.19\linewidth]{results/CIFAR_result/1_RDueling_ship.png} \\
            \scriptsize Original & \scriptsize PGD & \scriptsize RGD & \scriptsize ZO-RGD & \scriptsize RDNGD \\
        \end{tabular}
        \caption{Visual comparison of adversarial examples}
        \label{fig:cifar_visual}
    \end{subfigure}
    

    % --------------------------------------------------
    % Part (b) & (c): Loss Plots
    % --------------------------------------------------
    \begin{subfigure}[b]{0.48\linewidth}
        \centering
        \includegraphics[width=\linewidth,height=1.1in]{results/CIFAR_result/1_loss_vs_iteration.png}
        \caption{Adversarial loss vs. Iteration}
        \label{fig:cifar_loss_iter}
    \end{subfigure}
    \hfill 
    \begin{subfigure}[b]{0.48\linewidth}
        \centering
        \includegraphics[width=\linewidth,height=1.1in]{results/CIFAR_result/1_loss_vs_time.png}
        \caption{Adversarial loss vs. time}
        \label{fig:cifar_loss_time}
    \end{subfigure}

    \vspace{5pt}

    % --------------------------------------------------
    % Main Caption
    % --------------------------------------------------
    \caption{Attack results on CIFAR-10. (a) Generated adversarial images. (b) and (c) show the convergence curves. Our estimator yields accurate gradient estimation with only $10$ samples (vs. $500$ for ZO-RGD), demonstrating superior query efficiency.}
    \label{fig:cifar_attack_example}
\end{figure}



Figure~\ref{fig:cifar_attack_example} illustrates representative adversarial examples (additional results in Section~\ref{sec:additiona}) generated by different attacks. The white-box methods PGD and RGD are included as references but are infeasible in our comparison oracle setting. For all these test images, the attacks successfully induce misclassification with visually indistinguishable perturbations to the original image. Among the practical black-box methods, ZO-RGD and our RDNGD attack, the loss curves show that RDNGD yields a more stable optimization trajectory and reaches higher adversarial loss in fewer iterations and less CPU time, indicating a clear efficiency advantage. Moreover, as noted earlier, RDNGD operates with strictly weaker oracle information than ZO-RGD which makes its performance gains even more remarkable.


\subsubsection{Horizon leveling}
In this experiment, we evaluate our methods in the horizon leveling problem presented in Section~\ref{sec:example}. For each image in the HLW dataset~\citep{Workman2016}, we compute $R_{\text{tilt}} \in \mathrm{SO}(2)$ between the human-annotated horizon and the horizontal axis. Our goal is to find the correction rotation $R^\star \in \mathrm{SO}(2)$ that minimizes the misalignment cost $f(R)$:
\iffalse
{\color{red}In this experiment, we evaluate our methods in the horizon leveling problem presented in Section~\ref{sec:example}. For each image in the HLW dataset~\citep{Workman2016}, the human-annotated horizon is given by two endpoints in the image plane. From these endpoints we compute the tilt angle of the horizon relative to the horizontal axis, which we represent as a tilt rotation $R_{\text{tilt}} \in \mathrm{SO}(2)$. Our goal is to find the correction rotation $R^\star \in \mathrm{SO}(2)$ that minimizes the misalignment cost $f(R)$:}
\fi
\[
    \min_{R \in \mathrm{SO}(2)} f(R):= \| R R_{\text{tilt}} - I \|_F^2.
\]
While dueling feedback could be obtained from human comparisons, we use $f(R)$ as a scalable and reproducible surrogate for human preference. Given two rotations $R_1$ and $R_2$, the oracle returns the one with smaller misalignment, providing pairwise preferences on $\mathrm{SO}(2)$ without revealing function values.
We set the maximum iteration number $T=100$, $\nu=10^{-6}$, and constant step size $\eta=10^{-2}$.
\begin{figure}[t]
    \centering
    \begin{subfigure}[b]{0.31\linewidth}
        \centering
        \includegraphics[width=\linewidth]{results/Horizon_result/3744673249_79bcf131fc_o_original.png}
        \caption{Original}
    \end{subfigure}
    \hfill
    \begin{subfigure}[b]{0.31\linewidth}
        \centering
        \includegraphics[width=\linewidth]{results/Horizon_result/3744673249_79bcf131fc_o_correction.png}
        \caption{Corrected}
    \end{subfigure}
    \hfill
    \begin{subfigure}[b]{0.35\linewidth}
        \centering
        \includegraphics[width=\linewidth]{results/Horizon_result/3744673249_79bcf131fc_o_loss_vs_iteration.png}
        \caption{Loss vs Iteration}
    \end{subfigure}

    \caption{Horizon leveling results. The loss stabilizes at $10^{-5}$ due to a large constant step size, chosen to balance convergence speed and the optimality gap.}
    \label{fig:horizon_results_example}
\end{figure}

\iffalse
\begin{figure}[t]
    \centering
    \setlength{\tabcolsep}{4pt} 
    \renewcommand{\arraystretch}{0.5}
    
    \begin{tabular}{ccc}
        \includegraphics[width=0.26\linewidth]{results/Horizon_result/3744673249_79bcf131fc_o_original.png} &
        \includegraphics[width=0.26\linewidth]{results/Horizon_result/3744673249_79bcf131fc_o_correction.png} &
        \includegraphics[width=0.32\linewidth]{results/Horizon_result/3744673249_79bcf131fc_o_loss_vs_iteration.png}
    \end{tabular}

    \caption{Horizon leveling results. The loss stabilizes at $10^{-5}$ due to a large constant step size, chosen to balance convergence speed and the optimality gap.}
    
    \label{fig:horizon_results_example}
\end{figure}
\fi




Figure~\ref{fig:horizon_results_example} shows an example of our dueling horizon-leveling procedure on real image (additional results in Section~\ref{sec:additiona}). Across all cases, suboptimality reaches $\sim 10^{-5}$ within $30$ iterations, showing high accuracy with few iterations despite no access to function values or gradients. In each example, the corrected image looks visually level indicating that the algorithm successfully recovers a good rotation.
%, with the horizon or main structures no longer tilted, indicating that the algorithm successfully gets a visually good rotation.
\section{Conclusion}
In this paper, we establish the first theoretical framework for Riemannian dueling optimization that uses only comparison oracles. We proposed three algorithms for tackling Riemannian dueling optimization in different scenarios: (i) geodesically $L$-smooth objective; (ii) geodesically $L$-smooth and geodesically (strongly) convex objective where projection onto the constraint set $\mathcal{X}$ is cheap; (iii) geodisically $L$-smooth and geodesically convex objective when projection is prohibited. We established iteration and oracle complexities for each scenario. Moreover, our refined analysis overcomes intrinsic geometric barriers and improves upon existing Euclidean results by offering tighter constants and more flexible hyperparameter selection. Numerical experimental results on both synthetic problems and real applications demonstrated the capability of the proposed algorithms. Our work opens several promising directions on Riemannian optimization with dueling feedback, e.g., accelerated algorithms, Hessian-aware methods for identifying local minima, and exploring whether the noise in direction estimator can facilitate escape from saddle points.


\label{author info}
\label{final author}


\bibliography{reference}
\bibliographystyle{icml2026}
\newpage
\appendix
\onecolumn


\section{Preliminaries on Riemannian Geometry }\label{appendix:prelim}
Let $\mathcal{M}$ be a Riemannian manifold. We introduce some relevant geometric concepts and definitions below.
\begin{definition}[\bf Tangent Vector and Tangent Space]\label{def:tangent_space} 
Let $\mathcal{M}$ be a Riemannian manifold and let $x$ be a point on the manifold $\mathcal{M}$. Define $\mathfrak{F}_x(\mathcal{M})$ to be the set of smooth functions on $\mathcal{M}$ defined on the neighborhood of $x \in\mathcal{M}$. A tangent vector $v_x:\mathfrak{F}_x(\mathcal{M})\rightarrow\mathbb{R}$ at $x$ is a linear operator such that $$v_x(f)=\frac{d(f(\gamma(t)))}{dt}\biggr\rvert_{t=0}, f\in\mathfrak{F}(\mathcal{M}),$$where $\gamma:(-\delta,\delta)\rightarrow\mathcal{M}$ is a smooth curve on $\mathcal{M}$ with $\gamma(0)=x$. The tangent space $\mathrm{T}_x\mathcal{M}$ is the linear space of all tangent vector at $x$.
\end{definition}

\begin{definition}[\bf Riemannian Gradient]\label{def:riemannian_gradient} 
Suppose $f$ is a smooth function on $\mathcal{M}$. The Riemannian gradient $\mathrm{grad} f(x)$ is the unique tangent vector in $\mathrm{T}_x \mathcal{M}$ satisfying
\begin{align*}
\langle v_x, \mathrm{grad} f(x) \rangle_x = v_x(f)\quad\forall v_x \in \mathrm{T}_x \mathcal{M}.
\end{align*}
\end{definition}

Geodesics extend the conceppt of straight line to manifold, and represents the locally distance-minimizing paths.

\begin{definition}[\bf Geodesic]\label{def:geodesic} 
Let $ (\mathcal{M}, g) $ be a Riemannian manifold. A smooth curve $ \gamma: I \to \mathcal{M} $, where $ I \subset \mathbb{R} $ is an interval, is called a geodesic if it satisfies that
\begin{align*}
\nabla_{\dot{\gamma}(t)} \dot{\gamma}(t) = 0 \quad \text{for all } t \in I,
\end{align*}
where $ \nabla $ is the Levi-Civita connection associated with the metric $ g $.
\end{definition}



The exponential mapping at a point $ x \in M $ maps a tangent vector $ v \in \mathrm{T}_x M $ to the point reached by the unique geodesic starting at $ x $ with initial velocity $ v $, evaluated at time $ t = 1 $. 

\begin{definition}[\bf Exponential Mapping]\label{def:exponential_mapping} 
Let $ \mathcal{M} $ be a Riemannian manifold and $ x \in M $. The exponential mapping at $ x $, denoted $ \mathrm{Exp}_x: \mathrm{T}_x \mathcal{M} \to \mathcal{M} $, is defined by
\begin{align*}
\mathrm{Exp}_x(v) = \gamma_v(1),
\end{align*}
where $ \gamma_v: [0,1] \to \mathcal{M} $ is the unique geodesic such that $ \gamma_v(0) = x $ and $ \dot{\gamma}_v(0) = v \in \mathrm{T}_x \mathcal{M} $. That is, $ \mathrm{Exp}_x(v) $ is the endpoint of the geodesic starting at $ x $ with initial velocity $ v $, evaluated at time $ t = 1 $.
\end{definition}

\begin{definition}[\bf Riemannian distance]
Let $(\mathcal{M}, g)$ be a Riemannian manifold with metric tensor $g$. 
For any two points $x,y \in \mathcal{M}$, the Riemannian distance between $x$ and $y$ is defined as
\[
d(x,y) := \inf_{\gamma \in \Gamma(x,y)} \int_{0}^{1} 
    \sqrt{ g_{\gamma(t)}\bigl(\dot{\gamma}(t), \dot{\gamma}(t)\bigr)}  dt,
\]
where $\Gamma(x,y)$ denotes the set of all smooth curves $\gamma:[0,1]\to \mathcal{M}$ such that $\gamma(0)=x$ and $\gamma(1)=y$. 
\end{definition}


\begin{remark}
The exponential mapping $\mathrm{Exp}_x$ is a local diffeomorphism. That is, within a open neighborhood (restricted by the manifold), $\mathrm{Exp}_x$ is invertible and well defined. We use $\mathrm{Log}_{x}$ to denote the invert map $\mathrm{Exp}_x^{-1}$.
\end{remark}


Tangent vectors at two different points lie in different spaces, and thus cannot be directly compared. To solve this issue, parallel transport is defined to move a tangent vector along the geodesics in the meanwhile preserving the length and direction.
\begin{definition}[\bf Parallel Transport]\label{def:parallel_tranposrt} 
Let $ \gamma: [0,1] \to \mathcal{M} $ be a smooth curve on a Riemannian manifold $ \mathcal{M} $, and let $ V(t) \in T_{\gamma(t)} \mathcal{M} $ be a vector field along $ \gamma $. The vector field $ V $ is said to be parallel along $ \gamma $ if it satisfies
\begin{align*}
\nabla_{\dot{\gamma}(t)} V(t) = 0 \quad \text{for all } t \in [0,1].
\end{align*}
Given an initial vector $ v_0 \in T_{\gamma(0)} \mathcal{M} $, there exists a unique parallel vector field $ V(t) $ such that $ V(0) = v_0 $. The vector $ V(1) \in T_{\gamma(1)} \mathcal{M} $ is called the parallel transport of $ v_0 $ along $ \gamma $.
\end{definition}




\section{Some Useful Lemmas}
In the proof provided in the appendix, we adopt the following notation: for any positive integer $N$, let $[N]$ denote the set $\{0, 1, \ldots, N\}$. Additionally, we define $\mathcal{U}_k = \{u_0, u_1, \ldots, u_k\}$ as the history containing all information available up to step $k$.


\begin{lemma}\label{lem:cosine_annealing_step sizes}
Let $\{\eta_k\}_{k=0}^{T-1}$ be the cosine annealing step size defined in Definition \ref{def:cosine_annealing_step size}\footnote{Here, similar to \cite{li2021secondlookexponentialcosine}, we assume $\eta_{\min}=0$ for simplicity.}. It holds that:
\begin{itemize}[leftmargin=*]
  \item %\label{itm:cos-i} 
  $\displaystyle \sum_{k=0}^{T-1} \cos\left(\frac{\pi k}{T}\right) = 1.$
  \item %\label{itm:cos-ii} 
  $\displaystyle \sum_{k=0}^{T-1} \eta_k
      = \frac{\eta_0}{2}\left( T + \sum_{k=0}^{T-1}\cos\left(\frac{\pi k}{T}\right) \right)
      = \frac{\eta_0}{2}(T+1).$
  \item %\label{itm:cos-iii} 
  $\displaystyle \sum_{k=0}^{T-1} \eta_k^{2}
      = \frac{\eta_0^{2}}{8}(3T+4).$
\end{itemize}
\end{lemma}

\begin{proof}
We first prove part (i). According to \cite{li2021secondlookexponentialcosine}, it holds that $\sum_{k=1}^{T} \cos\left(\frac{k\pi}{T}\right) = -1$. Adjusting the summation range to $k=0, \dots, T-1$, we obtain
\begin{align*}
\sum_{k=0}^{T-1} \cos\left(\frac{k\pi}{T}\right) 
= \sum_{k=1}^{T} \cos\left(\frac{k\pi}{T}\right) - \cos(\pi) + \cos(0) 
= -1 - (-1) + 1 
= 1.
\end{align*}
Part (ii) follows immediately by substituting the result of part (i) into the definition of $\eta_k$:
\[
\sum_{k=0}^{T-1} \eta_k = \frac{\eta_0}{2}\left(T + \sum_{k=0}^{T-1}\cos\left(\frac{\pi k}{T}\right)\right) = \frac{\eta_0}{2}(T+1).
\]
For part (iii), expanding the square of $\eta_k$ yields
\begin{align*}
\sum_{k=0}^{T-1} \eta_k^{2}
&= \frac{\eta_0^{2}}{4}\sum_{k=0}^{T-1}\left(1 + 2\cos\left(\frac{\pi k}{T}\right) + \cos^{2}\left(\frac{\pi k}{T}\right)\right).
\end{align*}
Using the identity $\cos^2(x) = \frac{1+\cos(2x)}{2}$ and the fact that $\sum_{k=0}^{T-1} \cos(\frac{2\pi k}{T}) = 0$, we have
\[
\sum_{k=0}^{T-1} \cos^2\left(\frac{\pi k}{T}\right) = \sum_{k=0}^{T-1} \frac{1}{2} + \frac{1}{2}\sum_{k=0}^{T-1}\cos\left(\frac{2\pi k}{T}\right) = \frac{T}{2}.
\]
Substituting this result and the conclusion from part (i) back into the expansion gives
\begin{align*}
\sum_{k=0}^{T-1} \eta_k^{2}
&= \frac{\eta_0^{2}}{4}\left( T + 2(1) + \frac{T}{2} \right)
= \frac{\eta_0^{2}}{8}(3T+4).
\end{align*}
This completes the proof.
\end{proof}

\begin{lemma}\label{lem:riemannian_cosine_law}[\cite{zhang2016firstordermethodsgeodesicallyconvex} Corollary 8]
For any Riemannian manifold $\mathcal{M}$ where the sectional curvature is lower bounded by $\kappa$, the following holds for any $x, x_k \in \mathcal{X}$ and $x_{k+1} = \mathcal{P}_{\mathcal{X}}(\operatorname{Exp}_{x_k}(-\eta_k g_k))$:
\begin{equation*}
\langle -g_k, \Log_{x_k}(x) \rangle \leq \frac{1}{2\eta_k} (d^2(x_k, x) - d^2(x_{k+1}, x)) + \frac{\zeta(\kappa, d(x_k, x))\eta_k}{2} \|g_k\|^2, 
\end{equation*}
where $\zeta(\kappa, d(x_k, x)) = \frac{\sqrt{|\kappa|} d(x_k, x)}{\tanh(\sqrt{|\kappa|} d(x_k, x))}$.
\end{lemma}

\begin{lemma}\label{lem:xt_gt_upper_bound}
Suppose $f : \mathcal{M} \rightarrow \mathbb{R}$ is geodesically ${L}$-smooth. Consider the updates rule in Algorithm~\ref{alg:alg1}. Then we have for all $k \in [T]$,
\begin{align*}
d(x_{k}, x^*) \leq \sqrt{D} + \sum_{t=0}^{k-1}\eta_t 
\quad \text{and} \quad
\| \rgrad f(x_{k}) \| \leq {L} \left(\sqrt{D} + \sum_{t=0}^{k-1}\eta_t\right).
\end{align*}
\end{lemma}

\begin{proof}
Using non-expansiveness of the projection operator, we have that
\begin{align*}
d(x_{k} , x^* )
&\leq \sum_{t=0}^{k-1} d(x_{t+1} ,x_t ) + d( x_0 , x^* )\\
&= \sum_{t=0}^{k-1} d(\mathcal{P}_\mathcal{X}(\mathrm{Exp}_{x_t}(-\eta_t h_t)) ,\mathcal{P}_\mathcal{X}(x_t) ) + d( x_0 , x^* )\\
&\leq \sum_{t=0}^{k-1} d(\mathrm{Exp}_{x_t}(-\eta_t h_t) ,x_t ) + d( x_0 , x^* )\\
&\leq \sum_{t=0}^{k-1}\eta_t + \sqrt{D}.
\end{align*}
From that $f$ is geodesically ${L}$-smooth, we have
\begin{align*}
\| \rgrad f(x_{k})\|=
\| \rgrad f(x_{k}) - \Gamma_{x^*}^{x_{k}}\rgrad f(x^*) \|
\leq {L} d( x_{k} , x^* )
\leq {L} \left(\sqrt{D}+\sum_{t=0}^{k-1}\eta_t\right).
\end{align*}
This completes the proof. 
\end{proof}

\begin{lemma}\label{lem:gt_lower_bound}
Suppose $f : \mathcal{M} \rightarrow \mathbb{R}$ is a geodesically convex function. Then $f(x_k)-f(x^*)>\epsilon$ implies $\|\rgrad f(x_k)\|\geq {\epsilon}/{d(x_k,x^*)}.$
\end{lemma}

\begin{proof}
By geodesically convexity, we get 
\begin{align*}
    f(x^*)\geq f(x_k)+\langle \rgrad f(x_k), \Log_{x_k}(x^*)\rangle.
\end{align*}
Using the Cauchy-Schwarz inequality and that $d(x_k, x^*) = \|\Log_{x_k}(x^*)\|$, we obtain
\begin{align*}
    \epsilon&<f(x_k)-f(x^*)\leq - \langle \rgrad f(x_k), \Log_{x_k}(x^*)\rangle \leq\|\rgrad f(x_k)\|\cdot d(x_k,x^*).
\end{align*}
Dividing both sides by $d(x_k, x^*)$ completes the proof.
\end{proof}

\begin{lemma}\label{lem:inner_product_nt_logx}
Let $f:\mathcal{X}\subseteq\mathcal{M}:\rightarrow\mathbb{R}$ be a geodesically convex and geodesically ${L}$-smooth function, and $x^*$ be its minimizer over $\mathcal{X}$. For any $x_k\in\mathcal{X}$ such that $f(x_k)-f(x^*)\geq\epsilon$, we have $$\left\langle \frac{\rgrad f (x_k)}{\|\rgrad f(x_k)\|}, \frac{\Log_{x_k}(x^*)}{\|\Log_{x_k}(x^*)\|} \right\rangle \leq -\frac{1}{\|\Log_{x_k}(x^*)\|}\sqrt{\frac{ \epsilon}{2{L}}}.$$
\end{lemma}

\begin{proof}
From geodesically convexity, we have
$
f(x^*) \geq f(x_k) + \langle \rgrad f(x_k), \Log_{x_k}(x^*) \rangle_{x_k}.
$
Denote $g_k:=\rgrad f(x_k)$ and $w_k :=\frac{g_k}{\|g_k\|}$. 
This inequality can be rearranged as
\begin{align}\label{eq:convexity_bound}
\left\langle w_{k}, \Log_{x_k}(x^*)\right\rangle_{x_k} 
&\leq -\frac{f(x_k) - f(x^*)}{\|g_{k}\|} .
\end{align}
Next, from geodesically ${L}$-smoothness, for any tangent vector $u \in \mathrm{T}_{x_k}\mathcal{M}$, we have
\begin{align*}
f(\mathrm{Exp}_{x_k}(u)) \leq f(x_k) + \langle g_k, u \rangle_{x_k} + \frac{L}{2}\|u\|^2.
\end{align*}
Setting $u = -g_k/L$ and noting that $f(x^*) \leq f(\mathrm{Exp}_{x_k}(u))$, we get
$
f(x^*) \leq f(x_k) - \frac{1}{2L}\|g_k\|^2.
$
and thus 
$
\|g_{k}\| \leq \sqrt{2{L}\left(f(x_k) - f(x^*)\right)},
$
which, combining with \eqref{eq:convexity_bound} and $f(x_k) - f(x^*) \geq \epsilon$, yields
\begin{align*}
\left\langle w_{k}, \Log_{x_k}(x^*) \right\rangle_{x_k} 
\leq -\frac{f(x_k) - f(x^*)}{\sqrt{2{L}(f(x_k) - f(x^*))}} = -\sqrt{\frac{f(x_k) - f(x^*)}{2{L}}} \leq -\sqrt{\frac{ \epsilon}{2{L}}}.
\end{align*}
Finally, dividing both sides by $\|\Log_{x_k}(x^*)\|$ yields the desired result. 
\end{proof}

\begin{lemma}\label{lem:LG_smooth}
Let $f: \mathcal{X} \subseteq \mathcal{M} \to \mathbb{R}$ be a geodesically ${L}$-smooth function, and let $x^*$ be its minimizer over $\mathcal{X}$. For any $x \in \mathcal{X}$, we have 
$$\|\rgrad f(x)\|^2 \leq 2L(f(x) - f(x^*)).$$
\end{lemma}

\begin{proof}
By the definition of geodesic ${L}$-smoothness, for any tangent vector $u \in \mathrm{T}_{x}\mathcal{M}$, we have
\begin{align*}
    f(\mathrm{Exp}_{x}(u)) \leq f(x) + \langle \rgrad f(x), u \rangle_{x} + \frac{L}{2}\|u\|^2.
\end{align*}
Consider the update direction $u = -\frac{1}{L}\rgrad f(x)$. Plugging this into the inequality yields
\begin{align*}
    f(\mathrm{Exp}_{x}(u)) 
    &\leq f(x) - \frac{1}{L}\|\rgrad f(x)\|^2 + \frac{L}{2}\left\|-\frac{1}{L}\rgrad f(x)\right\|^2 \\
    &= f(x) - \frac{1}{2L}\|\rgrad f(x)\|^2.
\end{align*}
Since $x^*$ is the global minimizer, we have $f(x^*) \leq f(\mathrm{Exp}_{x}(u))$. Combining this with the upper bound derived above implies
$$f(x^*) \leq f(x) - \frac{1}{2L}\|\rgrad f(x)\|^2.$$
Rearranging the terms completes the proof.
\end{proof}



\section{Proof of Lemmas \ref{lem:grad_with_bias}, \ref{lem:estimating_normalized_gradient}, and Proposition \ref{prop:estimate_biased_NG}}
\subsection{Proof of Lemma\ref{lem:grad_with_bias}}\label{pf:grad_with_bias}
\begin{proof}
From Definition~\ref{as:lipgrad}, we have
\[
\begin{array}{rlcll}
\langle\pert u , \rgrad f(x)\rangle - \frac{L}{2}\pert^2 & \leq & f(\mathrm{Exp}_x(\pert u)) - f(x) & \leq & \langle\pert u , \rgrad f(x)\rangle + \frac{L}{2}\pert^2, \\ 
-\langle\pert u , \rgrad f(x)\rangle - \frac{L}{2}\pert^2 & \leq & f(\mathrm{Exp}_x(-\pert u))) - f(x) & \leq & -\langle\pert u , \rgrad f(x)\rangle + \frac{L}{2}\pert^2, 
\end{array}
\]
which further implies 
\begin{align*}
&2\langle\pert u , \rgrad f(x)\rangle - L\pert^2 \leq f(\mathrm{Exp}_x(\pert u))-f(\mathrm{Exp}_x(-\pert u))) \leq 2\langle\pert u , \rgrad f(x)\rangle +L\pert^2.%\numberthis\label{eq:fxplufxminusbound}
\end{align*}

Note that this indicates if $\abs{\langle u , \rgrad f(x)\rangle} >L\pert/2$, then
\begin{align}\label{lem:grad_with_bias-proof-sign-eq}
\sign (f(\mathrm{Exp}_x(\pert u)) - f(\mathrm{Exp}_x(-\pert u)))= \sign (\langle u , \rgrad f(x)\rangle).
\end{align}


We now establish an upper bound on $\mathbb{P}(|\langle u , \rgrad f(x)\rangle|\leq{L} \pert/2 )$. Denote $a=\frac{\rgrad f(x)}{\norm{\rgrad f(x)}}$, $b=\frac{{L}\pert}{2\norm{\rgrad f(x)}}$, and $Z = \langle u , a\rangle$. Therefore, $\mathbb{P}(|\langle u , \rgrad f(x)\rangle|\leq{L} \pert/2 ) = \mathbb{P}(|Z| \leq b)$. By the rotational symmetry of the sphere, we can assume ${a} = e_1$, which implies $Z = u_1$, the first coordinate of $u$. 

Similar to the proof of Lemma~\ref{pf:estimating_normalized_gradient}, the marginal density of $ Z \in [-1,1] $ is,
\begin{align*}
    f_Z(z) = A_d (1 - z^2)^{(d - 3)/2} \quad \text{for } z \in [-1,1], \mbox{ where } A_d = \frac{\Gamma\left(\frac{d}{2}\right)}{\sqrt{\pi}  \Gamma\left(\frac{d-1}{2}\right)}.
\end{align*}
Since $ (1 - z^2)^{(d - 3)/2} \leq 1 $ for all $ z \in [0,1] $, we have $\int_0^b f_Z(z) dz \leq A_d  b$, which further implies,
\begin{align*}
\mathbb{P}(|Z| \leq b) = \mathbb{P}(|\langle u , a\rangle| \leq b)=\int_{-b}^b f_Z(z) dz \leq 2 A_d  b \leq 2\sqrt{\frac{d}{2\pi}} b = \sqrt{\frac{d}{2\pi}}\frac{{L}\pert}{\norm{\rgrad f(x)}}, 
\end{align*}
where the second inequality is from the Gautschi’s inequality. Therefore, from \eqref{lem:grad_with_bias-proof-sign-eq} we have that  
\begin{align*}
     & \mathbb{P}(\sign (f(\mathrm{Exp}_x(\pert u)) - f(\mathrm{Exp}_x(-\pert u)))= \sign (\langle u , \rgrad f(x)\rangle)) \\
     \geq &  \mathbb{P}(|\langle u, \mathrm{grad} f(x) \rangle| > L \nu/2)
    \geq  1-\sqrt{\frac{d}{2\pi}}\frac{{L}\pert}{\norm{\rgrad f(x)}}.
\end{align*}
This completes the proof. 
\end{proof}
\subsection{Proof of Lemma \ref{lem:estimating_normalized_gradient} }\label{pf:estimating_normalized_gradient}
\begin{proof}
Note that, 
\begin{align*}
    \mathbb{E}[\operatorname{sign}(\langle v,u\rangle_x)u]=\mathbb{E}\left[\operatorname{sign}\left(\left\langle \frac{v}{\norm{v}_x},u\right\rangle_x\right)u\right].
\end{align*}
So, we can assume $\norm{v}_x=1$ without loss of generality. 
Let $P:\mathrm{T}_x\mathcal{M}\rightarrow \mathrm{T}_x\mathcal{M}$ denote the reflection operator along $v$ on $\mathrm{T}_x\mathcal{M}$, i.e., $P(\xi)= 2\langle v,\xi\rangle v-\xi$. For an arbitrary unit tangent vector $u\in\mathrm{T}_x\mathcal{M}$, denote $u':=P(u)$. We have that 
\begin{align*}
\operatorname{sign}( \langle v,u'\rangle) = \operatorname{sign} \left( 2 \|v\|^2 \langle v,u \rangle - \langle v,u \rangle \right) = \operatorname{sign}(\langle v,u\rangle).
\end{align*}
Since $u'\sim \mathrm{Unif}(\mathcal{S}_{\mathrm{T}_x\mathcal{M}}(1))$, we then have
\begin{align*}
\mathbb{E}[\operatorname{sign}(\langle v,u\rangle) u] 
&= \frac{1}{2} \mathbb{E}[\operatorname{sign}(\langle v,u\rangle) u] + \frac{1}{2} \mathbb{E}[\operatorname{sign}(\langle v,u'\rangle) u'] \\
&= \frac{1}{2} \mathbb{E}[\operatorname{sign}(\langle v,u\rangle) u] + \frac{1}{2} \mathbb{E}[\operatorname{sign}(\langle v,u\rangle)(2 (\langle v,u\rangle) v-u)] \\
&= \mathbb{E}[(\langle v,u\rangle) \operatorname{sign}(\langle v,u\rangle)] v\\
&=\mathbb{E}[|\langle v,u\rangle|]v\\
&=\pert v, %\numberthis \label{eq:vuinprod}
\end{align*}
where $\pert = \mathbb{E}[|\langle v,u\rangle|]$.
The remaining thing is to derive lower and upper bounds for $\pert$.


We first derive an upper bound for $\nu$. There exists an orthogonal matrix $R\in O(d)$ with $Rv=e_d$ where $e_d$ is a vector with all zeros but the last element is 1, and set $w:=Ru$. By symmetry of uniform distribution, $u'\coloneqq R^{-1}u$, and $w$ follows $\mathrm{Unif}(\mathcal{S}_{\mathrm{T}_x\mathcal{M}}(1))$ as well. Thus we have
$\mathbb{E}[|\langle v,u\rangle|]=\mathbb{E}[|\langle Rv,R^{-1}u\rangle|]=\mathbb{E}[|\langle e_1,u'\rangle|]=\mathbb{E}[|u'_1|]$.
We observe that 
\begin{align*}
\mathbb{E}[u_1^2] = \frac{1}{d} \mathbb{E}\left[\sum_{i=1}^d u_i^2\right] = \frac{1}{d},
\end{align*}
and thus get the following upper bound for $\pert$:
\begin{equation}\label{eq:nu-upper-bound}
\nu=\mathbb{E}[|\langle v,u\rangle|]=\mathbb{E}[|u'_1|]=\mathbb{E}[|u_1|] \leq \sqrt{\mathbb{E}[u_1^2]} = \frac{1}{\sqrt{d}}.
\end{equation}

Next we derive a lower bound for $\nu$. Using $\langle v,u\rangle_x = \langle Rv,Ru\rangle_x = \langle e_d, w\rangle_x$, we have
\begin{align}\label{eq:rotation}
\nu v = \mathbb{E}[\operatorname{sign}(\langle v,u\rangle_x)u]
=\mathbb{E}[\operatorname{sign}(\langle Rv,Ru\rangle_x)R^{\top}(Ru)]
=R^{\top}\mathbb{E}[\operatorname{sign}(\langle e_d,w\rangle_x)w].
\end{align}
Observe that the term $\operatorname{sign}(\langle e_d, w\rangle_x)$ depends only on the $d$-th coordinate $w_d$. Due to the rotational symmetry of the uniform distribution of $w$, for any $k \neq d$, the expectation of the $k$-th component $\mathbb{E}[\operatorname{sign}(w_d)w_k]$ vanishes. 
So $\mathbb{E}[\operatorname{sign}(\langle e_d,w\rangle_x)w]=\mathbb{E}[\operatorname{sign}(w_d)w]=C_d e_d$ for some scalar $C_d>0$.
Therefore, from \eqref{eq:rotation} we know that $\nu v = C_d R^\top e_d = C_d v$, which implies 
\begin{align*}
\nu=C_d=\langle \mathbb{E}[\operatorname{sign}(w_d)w],e_d\rangle_x=\mathbb{E}[\operatorname{sign}(w_d)w_d]=\mathbb{E}[|w_d|].
\end{align*}
Now, the marginal density of $ w_d \in [-1,1] $ is \citep{vershynin2009high}:
\begin{align*}
    f(w_d) = A_d (1 - w_d^2)^{(d - 3)/2}  \mbox{ where } A_d = \frac{\Gamma\left(\frac{d}{2}\right)}{\sqrt{\pi}  \Gamma\left(\frac{d-1}{2}\right)},
\end{align*}
which yields 
\begin{align}\label{nu-lower-bound}
    \nu=\mathbb{E}[|w_d|]=A_d\int_{-1}^1|t|(1 - t^2)^{(d - 3)/2}dt=2A_d\int_{0}^1t(1 - t^2)^{(d - 3)/2}dt
    \geq \frac{1}{\sqrt{2\pi d}},
\end{align}
where the last inequality follows the Gautschi’s inequality.

Note that $\hat{C}$ in \eqref{lem:estimating_normalized_gradient-eq} satisfies $\hat{C} = \nu \sqrt{d}$, and therefore, by combining \eqref{eq:nu-upper-bound} and \eqref{nu-lower-bound}, we obtain $\hat{C} \in [1/\sqrt{2\pi},1]$, and this completes the proof. 
\end{proof}





\subsection{Proof of Proposition \ref{prop:estimate_biased_NG}}
\begin{proof}\label{pf:estimate_biased_NG}
For any vector $v \in \mathcal{S}_{\mathrm{T}_x\mathcal{M}}(1)$, denote 
\[A:= \operatorname{sign}(f(\mathrm{Exp}_x( \pert u)) - f(\mathrm{Exp}_x( -\pert u)))  \langle u, v\rangle, \quad B:= \operatorname{sign}(\langle\rgrad f(x), u\rangle_x)  \langle u, v\rangle. \]
From Lemma \ref{lem:grad_with_bias} we know that $\mathbb{P}(A=B)\geq 1-\gamma_x$, which together with $|A-B|\leq 2|\langle u,v\rangle|\leq 2$, yields $|\mathbb{E}_u A-\mathbb{E}_u B|\leq 2\gamma_x$. From Lemma \ref{lem:estimating_normalized_gradient}, we know 
\begin{align*}
   \mathbb{E}_u B = \mathbb{E}_u\big[\operatorname{sign}(\langle \rgrad f(x),u\rangle_x)u\big]
= \frac{\hat{C}}{\sqrt{d}}\frac{\rgrad f(x)}{\|\rgrad f(x)\|},
\end{align*}
which completes the proof.
\end{proof}



\section{Proof of Theorem \ref{thm:beta_smoothness_complexityapp}}

\subsection{Proof of Theorem \ref{thm:beta_smoothness_complexityapp} Part (i)}\label{pf:beta_smoothness_complexity} 

\begin{remark}[Expectation Notation]
Unless specified otherwise by a subscript, the operator $\mathbb{E}[\cdot]$ denotes the total expectation taken over all random variables $\mathcal{U}_k$ generated up to the current iteration. We use the subscript notation (e.g., $\mathbb{E}_{u_k}[\cdot \mid x_k]$) for conditional expectations with respect to a specific random source $u_k$ given the history $x_k$.
\end{remark}

\begin{proof}
For the ease of notation, we denote $h_k=h_\pert(x_k)$ when there is no confusion. Choosing $v=\frac{\rgrad f(x_k)}{\|\rgrad f(x_k)\|}$ in Proposition \ref{prop:estimate_biased_NG}, we get 
\begin{equation}\label{eq:hkgradlower}
\begin{aligned}
\mathbb{E}_{u_k}\bigl[\langle h_k, \rgrad f(x_k)\rangle_{x_k}\mid x_k\bigr]
&\geq
\frac{\hat{C}}{\sqrt{d}}
\|\rgrad f(x_k)\|
-2\gamma_{x_k}\|\rgrad f(x_k)\|\\
&\geq 
\frac{\hat{C}}{\sqrt{d}}
\|\rgrad f(x_k)\|
-2\sqrt{\frac{d}{2\pi}}{L}\pert.
\end{aligned}
\end{equation}
Since $\mathcal{X}=\mathcal{M}$, we have $x_{k+1} = \mathcal{P}_{\mathcal{X}}\left(\mathrm{Exp}_{x_k}(- \eta h_\pert(x_k))\right) = \mathrm{Exp}_{x_k}(- \eta h_\pert(x_k)).$ 
By Assumption~\ref{as:lipgrad}, we get 
\begin{align*}
    f(x_{k+1})\leq f(x_k)-\eta\langle h_k,\rgrad f(x_k)\rangle_{x_k}+\frac{L\eta^2}{2},
\end{align*}
which implies 
\begin{equation}\label{eq:hkgradupper}
\mathbb{E}_{u_k}\bigl[\langle h_k, \rgrad f(x_k))\rangle_{x_k}\mid x_k\bigr]\leq \frac{1}{\eta }\mathbb{E}_{u_k}\bigl[f(x_k)-f(x_{k+1})\mid x_k\bigr] + \frac{\eta{L}}{2}.
\end{equation} 
Combining \eqref{eq:hkgradlower} and \eqref{eq:hkgradupper} yields
\begin{equation}\label{eq:combine_hkgrad_bounds}
    \|\rgrad f(x_k)\| \leq \frac{\sqrt{d}}{\hat{C}}\left(\frac{1}{\eta }\mathbb{E}_{u_k}\bigl[f(x_k)-f(x_{k+1})\,\big|\,x_k\bigr] + \frac{\eta{L}}{2}+2\sqrt{\frac{d}{2\pi}}{L}\pert\right).
\end{equation}
Summing over $k=0,1, \ldots, T-1$ and taking the total expectation on both sides, we apply the law of total expectation to obtain
\begin{equation}\label{eq:beta_smooth_inequality}
\begin{aligned}
    \mathbb{E}\left[\sum_{k=0}^{T-1}\|\rgrad f(x_k)\|\right] 
    &\leq \sum_{k=0}^{T-1}\frac{\sqrt{d}}{\hat{C}}\left(\frac{1}{\eta }\mathbb{E}\left[\mathbb{E}_{u_k}\bigl[f(x_k)-f(x_{k+1})\,\big|\,x_k\bigr]\right] + \frac{\eta{L}}{2}+2\sqrt{\frac{d}{2\pi}}{L}\pert\right)\\
    &= \frac{\sqrt{d}}{\eta \hat{C} }\sum_{k=0}^{T-1}\mathbb{E}\bigl[f(x_k)-f(x_{k+1})\bigr] +\frac{T\sqrt{d}}{\hat{C}}\frac{\eta{L}}{2}+\frac{Td}{\hat{C}}2\sqrt{\frac{1}{2\pi}}{L}\pert\\
    &= \frac{\sqrt{d}}{\eta \hat{C} }\mathbb{E}\bigl[f(x_0)-f(x_{T})\bigr] + \frac{T\sqrt{d}}{\hat{C}}\frac{\eta{L}}{2}+\frac{Td}{\hat{C}}2\sqrt{\frac{1}{2\pi}}{L}\pert\\
    &\leq \frac{\sqrt{d}}{\eta \hat{C} }(f(x_0)-f^*) + \frac{T\eta{L}\sqrt{d}}{2\hat{C}}+\frac{2T{L}\pert d}{\hat{C}\sqrt{2\pi}}.
\end{aligned}
\end{equation}
Let $R$ be a random variable uniformly distributed over $\{0, \dots, T-1\}$, i.e., $\mathbb{P}(R=k) = 1/T$. The expected gradient norm at $x_R$ is given by
\[
\mathbb{E}[\|\rgrad f(x_R)\|] = \frac{1}{T} \sum_{k=0}^{T-1} \mathbb{E}[\|\rgrad f(x_k)\|].
\]
Dividing \eqref{eq:beta_smooth_inequality} by $T$ and using $\eta = \frac{1}{\sqrt{T}}$, we have
\begin{align*}
    \mathbb{E}\left[\|\rgrad f(x_R)\| \right]
    &\leq \frac{1}{T}\frac{\sqrt{d}}{\eta \hat{C} }(f(x_0)-f^*) + \frac{\sqrt{d}}{\hat{C}}\frac{\eta{L}}{2}+\frac{d}{\hat{C}}2\sqrt{\frac{1}{2\pi}}{L}\pert\\
    &{\leq}  \frac{1}{T}\left(\frac{\sqrt{d}}{\eta \hat{C} }\frac{{L} D}{2}\right) + \frac{\sqrt{d}}{\hat{C}}\frac{\eta{L}}{2}+\frac{d}{\hat{C}}2\sqrt{\frac{1}{2\pi}}{L}\pert\\
    &= \frac{1}{\sqrt{T}}\frac{\sqrt{d}}{\hat{C}}\frac{{L}D}{2} + \frac{1}{\sqrt{T}}\frac{\sqrt{d}}{\hat{C}}\frac{{L}}{2}+\frac{d}{\hat{C}}2\sqrt{\frac{1}{2\pi}}{L}\pert\\
    &= \frac{1}{2\sqrt{T}}\left(\frac{{L}\sqrt{d}}{\hat{C}} (D+1)\right)+\frac{2dL}{\hat{C}\sqrt{2\pi}}\pert.
\end{align*}
Choosing $T$ and $\pert$ as specified in Theorem \ref{thm:beta_smoothness_complexityapp} Part (i) yields the desired result.
\end{proof}





\subsection{Proof of Theorem \ref{thm:beta_smoothness_complexityapp} Part (ii)}\label{pf:beta_smoothness_complexity_cosine}

\begin{proof}
The derivation is analogous to the previous proof, with the only difference being the choice of step size. By simply substituting the cosine annealing step size $\eta_k$ for $\eta$ in \eqref{eq:hkgradlower}, \eqref{eq:hkgradupper}, and \eqref{eq:combine_hkgrad_bounds}, and summing over $k=0, \dots, T-1$, the argument proceeds identically. Following similar argument as \eqref{eq:beta_smooth_inequality}, we have
\begin{align*}
\frac{\hat{C}}{\sqrt{d}}
\sum_{k=0}^{T-1} \eta_k 
\mathbb{E}\left[ \left\| {\rgrad} f(x_k) \right\| \right]
\le
 f(x_0) - f^* + \frac{{L}}{2} \sum_{k=0}^{T-1} \eta_k^{2}
+ 2\sqrt{\frac{d}{2\pi}}{L}\pert \sum_{k=0}^{T-1} \eta_k.
\end{align*}
We define a random variable $R$ taking values in $\{0, \dots, T-1\}$ with probability $\mathbb{P}(R=k) = \frac{\eta_k}{\sum_{j=0}^{T-1} \eta_j}$. Then, the left-hand side can be rewritten as
\[
\frac{\hat{C}}{\sqrt{d}} \left(\sum_{k=0}^{T-1} \eta_k\right) \mathbb{E}[\|\rgrad f(x_R)\|].
\]
Dividing both sides by $\frac{\hat{C}}{\sqrt{d}}\sum_{k=0}^{T-1} \eta_k$ and using Lemma~\ref{lem:cosine_annealing_step sizes}, we obtain
\begin{align*}
\mathbb{E}\left[ \|\rgrad f(x_R)\| \right]
&\leq
\frac{\sqrt{d}}{\hat{C}}\Bigg[
    \frac{{L} D}{\frac{\eta_0}{2}(T+1)}
    + \frac{{L}}{2} \cdot \frac{\frac{\eta_0^2}{8}(3T+4)}{\frac{\eta_0}{2}(T+1)}
    + 2\sqrt{\frac{d}{2\pi}}{L}\pert
\Bigg]\\
&= 
\frac{\sqrt{d}}{\hat{C}}\Bigg[
    \frac{2{L} D}{\eta_{0}(T+1)}
    + \frac{{L} \eta_{0}}{8}\cdot\frac{3T+4}{T+1}
    + 2\sqrt{\frac{d}{2\pi}}{L}\pert
\Bigg]\\
&\leq 
\frac{\sqrt{d}}{\hat{C}}\Bigg[
    \frac{2{L} D}{\eta_{0}(T+1)}
    + \frac{4{L} \eta_{0}}{8}
    + 2\sqrt{\frac{d}{2\pi}}{L}\pert
\Bigg]\\
&\leq 
\frac{\sqrt{d}}{\hat{C}}\Bigg[
    \frac{2{L} D}{\sqrt{T}}
    + \frac{{L}}{2\sqrt{T}}
    + 2\sqrt{\frac{d}{2\pi}}{L}\pert
\Bigg].
\end{align*}
Choosing $T$ and $\pert$ as in \eqref{thm:beta_smoothness_complexityapp-cosine-parameter} yields the desired result.
\end{proof}




\section{Proof of Theorem \ref{thm:beta_smoothness_convex_complexityapp}}

\subsection{Proof of Theorem \ref{thm:beta_smoothness_convex_complexityapp} Part (i)}\label{pf:beta_smoothness_convex_complexity}
\begin{proof}
We prove by contradiction. Note that the algorithm returns $\hat{x}_T$, and $f(\hat{x}_T) = \min_{0\le k\le T} f(x_k)$. Suppose that $f(\hat{x}_T) \geq f(x^*) + \epsilon$. This implies $f(x_k) \geq f(x^*) + \epsilon$ for all $k \in \{0, \dots, T-1\}$. 
We have $\bar\zeta\geq\zeta(\kappa,d(x_k,x^*))$ for all $x_k\in\mathcal{X}$, as $\zeta$ is increasing in the second entry and $\mathfrak{D}$ is the diameter of $\mathcal{X}$. 
By Lemma~\ref{lem:xt_gt_upper_bound}, we get $d(x_{k},x^*)\leq \sqrt{D}+\eta k \leq \sqrt{D}+\eta T$ for $k < T$. By Lemma~\ref{lem:gt_lower_bound}, we get $\|\rgrad f(x_k)\|\geq \frac{\epsilon}{d(x_k,x^*)}$. Thus, $\nu$ defined in \eqref{thm:beta_smoothness_convex_complexityapp-constant-step-param} satisfies
\begin{align}\label{nu-upper-bound}
\pert = \frac{1}{4\sqrt{2}d} \frac{(\epsilon/L)^{3/2}}{(\sqrt{D} + \eta T)^2} \leq \frac{1}{4\sqrt{2}d} \frac{(\epsilon/L)^{3/2}}{d(x_k, x^*)^2} \leq \frac{1}{4\sqrt{2}Ld} \frac{\|\rgrad f(x_k)\|}{d(x_k, x^*)} \sqrt{\frac{\epsilon}{L}}, \forall k\in[T].
\end{align}
From Lemma~\ref{lem:riemannian_cosine_law} and Proposition~\ref{prop:estimate_biased_NG}, we have the recursive bound:
\begin{align}
\mathbb{E}_{u_k}\left[d^2(x_{k+1}, x^*)|x_k\right] 
&\leq d^2(x_k ,x^*) -\frac{1}{\sqrt{\pi}}\frac{\sqrt{\epsilon}}{\sqrt{d{L}}}\eta + {\left(\sqrt{\frac{8d}{\pi}}\frac{{L}\|\Log_{x_k}(x^*)\|}{\norm{\rgrad f(x_k)}}\pert\right)}\eta + \bar\zeta\eta^2\nonumber\\
&\leq d^2(x_k ,x^*) -\frac{1}{\sqrt{\pi}}\frac{\sqrt{\epsilon}}{\sqrt{d{L}}}\eta +\frac{1}{2\sqrt{\pi}}\frac{\sqrt{\epsilon}}{\sqrt{d{L}}}\eta + \bar\zeta\eta^2\nonumber\\
&=  d^2(x_k ,x^*) - \frac{\epsilon}{16\pi d{L}\bar{\zeta}},\label{eq:beta_smooth_g_convex}
\end{align}
where the equality is due to $\eta$ defined in \eqref{thm:beta_smoothness_convex_complexityapp-constant-step-param}. 

Taking the full expectation on both sides of \eqref{eq:beta_smooth_g_convex} and iterating the inequality recursively for $k=0, \ldots, T-1$, we derive 
\begin{equation*}\label{eq:dist_ub}
\begin{aligned}
0\leq \mathbb{E}\left[ d^2(x_{T}, x^*) \right]
\leq d^2(x_0 ,x^*) - \frac{\epsilon}{16\pi d{L}\bar{\zeta}}T < 0,
\end{aligned}
\end{equation*}
where the last inequality is due to the choice of $T$ in \eqref{thm:beta_smoothness_convex_complexityapp-constant-step-param}. This is a contradiction. 
Thus, there must exist an iteration $k$ such that $f(x_k) - f(x^*) < \epsilon$, which implies $f(\hat{x}_T) - f(x^*) < \epsilon$.
\end{proof}



\subsection{Proof of Theorem \ref{thm:beta_smoothness_convex_complexityapp} Part (ii)}\label{pf:beta_smoothness_convex_complexity_cosine}
\begin{proof}
We follow the same contradiction argument as in Part (i). Suppose that $f(x_k) \geq f(x^*) + \epsilon$ for all $k \in \{0, \dots, T-1\}$.
Using the cosine annealing step size $\eta_k$, the accumulated step length is $\sum_{j=0}^{T-1}\eta_j = \frac{\eta_0}{2}(T+1)$. For the choice of $\pert$ in \eqref{thm:beta_smoothness_convex_complexityapp-cosine-step-param}, we use the bound involving $T$ as a sufficient approximation for large $T$, ensuring \eqref{nu-upper-bound} holds.

Replacing $\eta$ with $\eta_k$ in the derivation of \eqref{eq:beta_smooth_g_convex}, we obtain:
\begin{align*}
\mathbb{E}_{u_k}\left[d^2(x_{k+1}, x^*)|x_k\right] 
&\leq d^2(x_k ,x^*) -\frac{1}{\sqrt{\pi}}\frac{\sqrt{\epsilon}}{\sqrt{d{L}}}\eta_k + \frac{1}{2\sqrt{\pi}}\frac{\sqrt{\epsilon}}{\sqrt{d{L}}}\eta_k + \bar\zeta\eta_k^2 \\
&= d^2(x_k ,x^*) - \frac{\sqrt{\epsilon}}{2\sqrt{\pi d{L}}}\eta_k + \bar\zeta\eta_k^2.
\end{align*}
Taking the full expectation and summing the inequality over $k=0, \dots, T-1$, we get
\begin{align*}
\mathbb{E}\left[ d^2(x_{T}, x^*) \right]
&\leq d^2(x_0, x^*) - \frac{\sqrt{\epsilon}}{2\sqrt{\pi d{L}}} \sum_{k=0}^{T-1} \eta_k + \bar\zeta \sum_{k=0}^{T-1} \eta_k^2.
\end{align*}
By Lemma~\ref{lem:cosine_annealing_step sizes}, we have $\sum_{k=0}^{T-1} \eta_k = \frac{\eta_0}{2}(T+1)$ and $\sum_{k=0}^{T-1} \eta_k^2 = \frac{\eta_0^2}{8}(3T+4)$.
Using the bound $\sum_{k=0}^{T-1} \eta_k^2 \leq \eta_0 \sum_{k=0}^{T-1} \eta_k$ (since $\frac{3T+4}{T+1} \le 4$), we have:
\begin{align*}
\mathbb{E}\left[ d^2(x_{T}, x^*) \right]
&\leq D - \left( \frac{\sqrt{\epsilon}}{2\sqrt{\pi d{L}}} - \bar\zeta \eta_0 \right) \sum_{k=0}^{T-1} \eta_k \\
&= D - \frac{\sqrt{\epsilon}}{4\sqrt{\pi d{L}}} \sum_{k=0}^{T-1} \eta_k \\
&= D - \frac{\sqrt{\epsilon}}{4\sqrt{\pi d{L}}} \frac{\eta_0}{2}(T+1).
\end{align*}
Substituting $\eta_0 = \frac{\sqrt{\epsilon}}{4\sqrt{\pi}\bar{\zeta}\sqrt{d{L}}}$:
\begin{align*}
\mathbb{E}\left[ d^2(x_{T}, x^*) \right]
&< D - \frac{\sqrt{\epsilon}}{8\sqrt{\pi d{L}}} \left( \frac{\sqrt{\epsilon}}{4\sqrt{\pi}\bar{\zeta}\sqrt{d{L}}} \right) T \\
&= D - \frac{\epsilon}{32\pi d{L}\bar{\zeta}}T.
\end{align*}
Given $T \geq 1 + \frac{32\pi d{L}\bar{\zeta}}{\epsilon}D$, we have $\mathbb{E}\left[ d^2(x_{T}, x^*) \right] < 0$, which is a contradiction.
\end{proof}




\section{Proof of Theorem \ref{thm:strong_convexity_complexity}}\label{pf:strong_convexity_complexity}

\begin{proof}
Note that in each iteration of Algorithm \ref{alg:alg2}, we run Algorithm \ref{alg:alg1} for $t = \lceil{1+64\pi d L\bar{\zeta}/\alpha}\rceil$ iterations. 
From Theorem~\ref{thm:beta_smoothness_convex_complexityapp} (i), with the choices of $\eta_k$ and $\nu_k$ defined in Theorem~\ref{thm:strong_convexity_complexity}, we have that for any phase $k \in \{0, \dots, K-1\}$, starting from $x^{(k)}$ with initial distance bound $D_{k} \geq\mathbb{E}[d^2(x^{(k)}, x^*)]$:
\begin{align*}
\mathbb{E}[f({x}^{(k+1)})] - f(x^*) \leq \epsilon_k = \frac{\alpha D_{k}}{4},
\end{align*}
with oracle complexity $ 2t $. Since $f$ is geodesically $\alpha$-strongly convex, we have
\begin{align*}
\mathbb{E}\left[\frac{\alpha}{2} d^2({x}^{(k+1)}, x^*)\right] \leq \mathbb{E}[f(x^{(k+1)})] - f(x^*) \leq \frac{\alpha D_{k}}{4},
\end{align*}
which implies 
\begin{align*}
\mathbb{E}[d^2({x}^{(k+1)}, x^*)] \leq \frac{D_{k}}{2} = D_{k+1}.
\end{align*}
This justifies the update rule in the algorithm. 
After $K$ phases, the expected distance squared of the final output $x^{(K)}$ is bounded by:
\begin{align*}
    \mathbb{E}[d^2(x^{(K)}, x^*)] \leq \frac{D}{2^K}.
\end{align*}

Note that Assumption \ref{as:lipgrad} and the fact that $\mathcal{X}$ is a bounded set imply that $f$ is geodesically Lipschitz continuous with some Lipschitz constant $l$. This further implies
\begin{align*}
    \mathbb{E}[f(x^{(K)}) - f(x^*)] 
    \leq l \cdot \mathbb{E}[d(x^{(K)} , x^*)] \leq l \sqrt{\mathbb{E}[d^2(x^{(K)} , x^*)]} \leq l \sqrt{\frac{D}{2^K}} \leq \epsilon,
\end{align*}
where the second inequality is Jensen's inequality, and the last inequality is due to the definition of $K$.
\end{proof}

\begin{remark}
All results developed in Section \ref{sec:section3} use the search direction $h_{\pert}$ defined in (\ref{eq:graddirestimator}), which is based on a single pair of samples. It is worth noting that all results in Section \ref{sec:section3} also hold for the batch averaged estimator $\bar{h}$ defined in (\ref{eq:batchgrad}). This is because our analysis relies solely on the expectation of $h_{\pert}$, which is identical to the expectation of $\bar{h}$. For more details, please refer to (\ref{eq:same_expectation}).
\end{remark}




\section{Proof of Theorem \ref{thm:RFW_batch}}\label{pf:RFW_batch}
\begin{proof}
We denote $\mathfrak{M}_f:={L}\mathfrak{D}^2$, $\tilde{C}=\sqrt{d}/\hat{C}$, normalized Riemannian gradient $w_k:= \frac{\operatorname{grad}f(x_k)}{\norm{\operatorname{grad}f(x_k)}}$, the function value gap by $\Delta_{k}:=f(x_{k}) - f(x^*)$ for $k=0,1,\ldots,T$. Let $\mathcal{B}_k := \{u^{(k)}_{1}, \dots, u^{(k)}_{M_k}\}$ denote the batch of i.i.d. random directions sampled at iteration $k$.

By geodesically $L$-smoothness, we have
\begin{align*}
\begin{aligned}
& \Delta_{k+1} = f(x_{k+1}) - f(x^*) \\
\leq & f(x_k) - f(x^*) 
   + s_k \norm{\operatorname{grad}f(x_k)}\langle w_k,  \Log_{x_k}(z_k) \rangle 
   + \tfrac{1}{2}s_k^2 \mathfrak{M}_f \\
   = & \Delta_{k}
   + s_k \norm{\operatorname{grad}f(x_k)}\langle \tilde{C}\bar h_k,  \Log_{x_k}(z_k) \rangle +s_k \norm{\operatorname{grad}f(x_k)}\langle w_k-\tilde{C}\bar h_k,  \Log_{x_k}(z_k) \rangle 
   + \tfrac{1}{2}s_k^2 \mathfrak{M}_f \\
   \leq & \Delta_{k}
   + s_k \norm{\operatorname{grad}f(x_k)}{\langle\tilde{C} \bar h_k,  \Log_{x_k}(x^*)\rangle} +s_k \norm{\operatorname{grad}f(x_k)}\langle w_k-\tilde{C}\bar h_k,  \Log_{x_k}(z_k) \rangle +
   \tfrac{1}{2}s_k^2 \mathfrak{M}_f \\
   = & \Delta_{k} 
   + s_k \norm{\operatorname{grad}f(x_k)}\langle w_k,  \Log_{x_k}(x^*) \rangle+ \tfrac{1}{2}s_k^2 \mathfrak{M}_f  +s_k \norm{\operatorname{grad}f(x_k)}\langle w_k-\tilde{C}\bar h_k,  \Log_{x_k}(z_k)-\Log_{x_k}(x^*)\rangle,
\end{aligned}
\end{align*}
where the second inequality is due to the definition of $z_k$ in Algorithm \ref{alg:RFW_batch}.
Taking expectation on both sides, we get
\begin{align}
&\mathbb{E}_{\mathcal{B}_k}\left[\Delta_{k+1}\mid x_k\right] \nonumber\\
\leq & \Delta_{k}
   + s_k \langle \mathrm{grad}f(x_k),  \Log_{x_k}(x^*) \rangle+ \tfrac{1}{2}s_k^2 \mathfrak{M}_f +s_k\norm{\operatorname{grad}f(x_k)}\mathbb{E}_{\mathcal{B}_k}[\langle w_k-\tilde{C}\bar h_k,\Log_{x_k}(z_k)-\Log_{x_k}(x^*)\rangle \mid x_k]\nonumber \\
   \leq & \Delta_{k}
   - s_k\Delta_k + \tfrac{1}{2}s_k^2 \mathfrak{M}_f +2\mathfrak{D}s_k \norm{\operatorname{grad}f(x_k)} \cdot\mathbb{E}_{\mathcal{B}_k}[\|w_k-\tilde{C}\bar h_k\|\mid x_k],\label{eq:batch_FW_eq}
\end{align}
where the last inequality used the geodesic convexity property \eqref{geo-convex-inequality}. 
Next we derive an upper bound for $\mathbb{E}_{\mathcal{B}_k}[\|{w_k-\tilde{C}\bar h_k}\|\mid x_k]$. We first show that the expectation of the batch-average estimator $\bar h_k$ is the same as the expectation of the single sample estimator $h_k:=h_{\pert_k}(x_k)$ defined in \eqref{eq:graddirestimator}. Denoting $\bar{m}_k:=\mathbb{E}_{\mathcal{B}_k}[\bar{h}_k\mid x_k ]$ and $m_k:=\mathbb{E}_{u_k}[h_k\mid x_k ]$, we have
\begin{equation}\label{eq:same_expectation}
\bar{m}_k=\mathbb{E}_{\mathcal{B}_k}[\bar{h}_k\mid x_k ]= \frac{1}{M_k} \sum_{j=1}^{M_k} \mathbb{E}_{u^{(k)}_j}[o^{(k)}_j u^{(k)}_j \mid x_k]= \mathbb{E}_{u_k}[o_k u_k\mid x_k ]= \mathbb{E}_{u_k}[h_k \mid x_k] = m_k.
\end{equation}
Next we bound the variance. By independence of $\{u_j\}_j$ and $\mathbb{E}_{u^{(k)}_j}[o^{(k)}_ju^{(k)}_j|x_k]-m_k=0$, we have 
\begin{align*}
    &\mathbb{E}_{\mathcal{B}_k}{[\|\bar{h}_k-\bar{m}_k\|_2^2\mid x_k]}=\frac{1}{M_k^2}\sum_{j=1}^{M_k}\mathbb{E}_{u^{(k)}_j}{[\|o^{(k)}_ju^{(k)}_j-m_k\|_2^2\mid x_k]}\\
    =&\frac{1}{M_k^2}\sum_{j=1}^{M_k}(\mathbb{E}_{u^{(k)}_j}{[\|o^{(k)}_ju^{(k)}_j\|_2^2\mid x_k]}-\|m_k\|_2^2)= \frac{1}{M_k}(1-\norm{m_k}^2).
\end{align*}
By Jensen’s inequality, we get
\begin{align}\label{thm8.4-bound-variance}
\mathbb{E}_{\mathcal{B}_k}\left[\|\bar h_k-\bar m_k\|\mid x_k\right]
\le \sqrt{\mathbb{E}_{\mathcal{B}_k}\left[\|\bar h_k-\bar m_k\|^2 \mid x_k\right]}
= \sqrt{{(1-\|m_k\|^2)}/{M_k}}
\le {1}/{\sqrt{M_k}}.
\end{align}
As a consequence of the above inequalities, we get that
\begin{equation}
\label{eq:batch_mean_plus_var}
\begin{aligned}
\mathbb{E}_{\mathcal{B}_k}[\| w_k-\tilde{C}\bar h_k \|\mid x_k]
&=\mathbb{E}_{\mathcal{B}_k}[\|(w_k-\tilde{C}\bar m_k)
     +\tilde{C}(\bar m_k-\bar h_k)\|\mid x_k]\\
&\le \| w_k-\tilde{C}\bar m_k\|
   +\tilde{C}\mathbb{E}_{\mathcal{B}_k}[\|\bar h_k-\bar m_k\|\mid x_k] \\
&= \tilde{C}\|w_k/\tilde{C}-\bar m_k\|
   +\tilde{C}\mathbb{E}_{\mathcal{B}_k}[\|\bar h_k-\bar m_k\|\mid x_k]\\
&\le 2\tilde{C}\gamma_{x_k}+\tilde{C}/\sqrt{M_k},
\end{aligned}
\end{equation}
where the last inequality is from \eqref{thm8.4-bound-variance} and Proposition \ref{prop:estimate_biased_NG} by setting $v = \frac{\bar{m}_k-w_k/\tilde{C}}{\|\bar{m}_k-w_k/\tilde{C}\|}$.

Combining \eqref{eq:batch_FW_eq} and \eqref{eq:batch_mean_plus_var}, we get
\begin{align*}
&\mathbb{E}_{\mathcal{B}_k}\left[\Delta_{k+1}\bigl|x_k\right]\\
 \leq & \Delta_{k} 
   - s_k \Delta_k + \tfrac{1}{2}s_k^2 \mathfrak{M}_f +2\mathfrak{D}s_k \norm{\operatorname{grad}f(x_k)} (2\tilde{C} {\gamma_{x_k}}+\tilde{C}/\sqrt{M_k})\\
 = & \Delta_{k} 
   - s_k \Delta_k+ \tfrac{1}{2}s_k^2 \mathfrak{M}_f +2\mathfrak{D}s_k \norm{\operatorname{grad}f(x_k)} \left(\tilde{C}\sqrt{\frac{2d}{\pi}}\frac{{L}\pert_k}{\|\rgrad f(x_k)\|}+\frac{\tilde{C}}{\sqrt{M_k}}\right)\\
 \leq & \Delta_{k} 
   - s_k \Delta_k+ \tfrac{1}{2}s_k^2 \mathfrak{M}_f +2\mathfrak{D}s_k\left(\tilde{C}\sqrt{\frac{2d}{\pi}}{L}\pert_k+ \frac{\sqrt{2{L}\Delta_k}\tilde{C}}{\sqrt{M_k}}\right)\\
 \leq & \Delta_{k} 
   - s_k \Delta_k+ \tfrac{1}{2}s_k^2 \mathfrak{M}_f +2\mathfrak{D}s_k\tilde{C}\sqrt{\frac{2d}{\pi}}{L}\pert_k+ \left(\frac{\Delta_k}{4} + \frac{8{L}\mathfrak{D}^2 d}{\hat{C}^2M_k}\right) s_k \\
   = &\left(1-\frac{3s_k}{4}\right)\Delta_k+ \frac{s_k^2}{2} (\mathfrak{M}_f+3),
\end{align*}
where the second inequality is from Lemma \ref{lem:LG_smooth}, the third inequality used Young's inequality, and the last equality is obtained by substituting the choices of $s_k$ in Algorithm \ref{alg:RFW_batch} and $\nu_k$ and $M_k$ in Theorem \ref{thm:RFW_batch}. 
Denote $a_k := 1-\frac{3}{4}s_k = 1-\frac{3/2}{k+3}$. Taking the total expectation on both sides of the above inequality and applying it recursively, we obtain %\eqref{eq:RFW_iteration_inequality} and telescoping the recursion yields
\begin{equation}\label{eq:RFW_inequality}
\mathbb{E}\left[\Delta_{k}\right] \leq 
{\left(\prod_{t=0}^{k-1}a_t\right) \Delta_0}
+ {\sum_{j=0}^{k-1} \frac{2(\mathfrak{M}_f+3)}{(j+3)^2}\prod_{t=j+1}^{k-1} a_t}.
\end{equation}
We now develop upper bounds for the two terms on the right hand side of \eqref{eq:RFW_inequality}. 
Note that 
\[
a_k = 1-\frac{3}{2}\cdot\frac{1}{k+3}\leq\left(1-\frac{1}{k+3}\right)^{3/2}=\left(\frac{k+2}{k+3}\right)^{3/2}.
\]
Therefore, 
\begin{equation*}
\prod_{t=0}^{k-1}a_t\leq\left(\prod_{t=0}^{k-1}\frac{t+2}{t+3}\right)^{3/2}=\left(\frac{2}{k+2}\right)^{3/2}\leq \frac{2}{k+2}, \mbox{ and } \prod_{t=j+1}^{k-1} a_t \leq  \left( \prod_{t=j+1}^{k-1}\frac{t+2}{t+3} \right)^{3/2}=  \left( \frac{j+3}{k+2}\right)^{3/2},
\end{equation*}
which further implies 
\begin{equation*}%\label{eq:RFW(ii)}
%\begin{aligned}
\sum_{j=0}^{k-1} \frac{2(\mathfrak{M}_f+3)}{(j+3)^2}\prod_{t=j+1}^{k-1} a_t\leq 
\sum_{j=0}^{k-1} 
   \frac{2(\mathfrak{M}_f+3)}{(j+3)^2} 
   \left(\frac{j+3}{k+2}\right)^{\tfrac{3}{2}}= \frac{2(\mathfrak{M}_f+3)}{(k+2)^{3/2}} 
   \sum_{j=0}^{k-1} \frac{1}{\sqrt{j+3}} \leq \frac{4(\mathfrak{M}_f+3)}{k+2},
%\end{aligned}
\end{equation*}
where the last inequality is from that $\sum_{m=3}^{k+2} \frac{1}{\sqrt{m}}
\leq \int_{2}^{k+2} x^{-1/2} dx \leq 2\sqrt{k+2}$.

Substituting these upper bounds to \eqref{eq:RFW_inequality}, we obtain
\begin{equation*}
\begin{aligned}
\mathbb{E}[\Delta_{k}]
\leq
\frac{2\Delta_{0} + 4\mathfrak{M}_f+12}{k+2}.  
\end{aligned}
\end{equation*}
So it takes $T=\lceil \frac{2\Delta_{0} + 4\mathfrak{M}_f+12}{\epsilon}\rceil$ to achieve
$\mathbb{E}[\Delta_{T}]\leq\epsilon$.
\end{proof}

\section{Sensitivity of Frank-Wolfe Search Direction}\label{sec:fw_sensitivity_example}
A key difference between projected gradient methods and Frank--Wolfe type methods is how they respond to noise in the gradient direction estimate. This effect is visible even in Euclidean case. In projected normalized gradient descent, the update takes the following form with noisy gradient direction $\tilde g_t$:
\(
     x_{t+1} = \mathcal{P}_{\mathcal{X}}(x_t - \eta \tilde g_t),
\)
Due to non-expansiveness (Assumption~\ref{as:proj}) of the projection operator, the distance between $x_{t+1}$, and the hypothetical update that one would have gotten with the true normalized gradient $g_t$, is upper bounded by $\eta\|\tilde g_t - g_t\|)$.

In contrast, Frank--Wolfe computes a search point via a linear minimization oracle
\[
     s(\tilde g_t) \in \arg\min_{s\in\mathcal{X}} \langle \tilde g_t, s-x\rangle,
\]
where $\mathcal{X}$ is a compact convex set. 
For such problems, the solution $s(\tilde g_t)$ lies on the boundary of $\mathcal{X}$. As a result, even an arbitrarily small perturbation of the direction $\tilde g_t$ can cause the oracle to select a completely different boundary point. In other words, even a small difference between $g_t$ and $\tilde g_t$, can lead to completely different search directions $s(g_t)$ and $s(\tilde g_t)$ and $\|s(g_t)-s(\tilde g_t)\|$ can be in the order of diameter $\mathcal{O}(\mathfrak{D})$ (see Figure~\ref{fig:fw_noise_diameter}). This means the variance of $s(\tilde g_t)$ is of constant order. Therefore, to guarantee descent and convergence for Frank--Wolfe with noisy or comparison-based oracles, we must explicitly control the variance of the direction estimator, whereas projected gradient methods are much more stable under the same level of noise. %Figure~\ref{fig:fw_noise_diameter} gives a simple numerical example of this noise amplifying effect.

% \begin{figure}[htbp]
%      \centering
%      \includegraphics[width=0.9\textwidth]{results/Frank_Wolfe_Noise_Diameter_Contrast.png}
%      \caption{Impact of the constraint set diameter on the sensitivity of the Frank-Wolfe Linear Minimization Oracle (LMO) to gradient noise. 
%      Left: A larger diameter results in a significant deviation between the LMO solutions computed from the true gradient versus a noisy estimate. 
%      Right: A smaller diameter reduces this deviation, even when the noise magnitude remains constant. 
%      This demonstrates how gradient noise error is amplified by the diameter of the constraint set.}
%      \label{fig:fw_noise_diameter}
% \end{figure}
% \clearpage
\begin{figure}[htbp]
    \centering
    \begin{subfigure}[b]{0.48\textwidth}
        \centering
        \includegraphics[width=\linewidth]{results/large_diameter.png} 
        \caption{Large diameter: Significant deviation between LMO solutions.}
        \label{fig:fw_large}
    \end{subfigure}
    \hfill
    \begin{subfigure}[b]{0.48\textwidth}
        \centering
        \includegraphics[width=\linewidth]{results/small_diameter.png
        }
        \caption{Small diameter: Reduced deviation with constant noise.}
        \label{fig:fw_small}
    \end{subfigure}
    
    \caption{Impact of the constraint set diameter on the sensitivity of the Frank-Wolfe Linear Minimization Oracle (LMO) to gradient noise. 
    (a) A larger diameter leads to significant deviation between the true and noisy LMO solutions. 
    (b) A smaller diameter mitigates this deviation, even when the noise magnitude remains constant. 
    Overall, this demonstrates how gradient noise error is amplified by the diameter of the constraint set.}
    \label{fig:fw_noise_diameter}
\end{figure}
\clearpage

\section{Additional Numerical Experiments}\label{sec:additiona}
\subsection{Additional Synthetic problem results}
We provide the CPU time results for the Rayleigh quotient problem in Figure~\ref{fig:rayleigh_cpu} and the Karcher Mean problem in Figure~\ref{fig:Karcher_mean_unconstrained_cpu}. In these low-dimensional synthetic settings, it can be observed that RDNGD converges more slowly than ZO-RGD. This occurs because RDNGD relies solely on a comparison oracle and lacks access to gradient magnitude information. Consequently, with an unit length search direction, RDNGD requires conservative step sizes to avoid oscillation around the minimum, whereas ZO-RGD benefits from gradient magnitude estimation for more efficient updates. This limitation is inherent to optimization problems utilizing only comparison oracles. However, for the high-dimensional real-world applications shown in Figure~\ref{fig:cifar_attack_example} and Figures~\ref{fig:cifar_ex1}--\ref{fig:cifar_ex3}, we observe a clear advantage in CPU time for RDNGD, demonstrating the efficiency of our method.

\begin{figure}[htbp]
    \centering
    % --- Rayleigh Quotient Figure ---
    \begin{subfigure}[b]{0.49\linewidth}
        \centering
        \includegraphics[width=\linewidth]{results/rayleigh_cpu_d=100.png}
        \caption{d=100}
        \label{fig:rayleigh_d100}
    \end{subfigure}
    \hfill
    \begin{subfigure}[b]{0.49\linewidth}
        \centering
        \includegraphics[width=\linewidth]{results/rayleigh_cpu_d=150.png}
        \caption{d=150}
        \label{fig:rayleigh_d150}
    \end{subfigure}
    
    \caption{CPUtime results for Rayleigh quotient maximization.}
    \label{fig:rayleigh_cpu}
\end{figure}

\begin{figure}[htbp]
    \centering
    % --- Karcher Mean Figure ---
    \begin{subfigure}[b]{0.49\linewidth}
        \centering
        \includegraphics[width=\linewidth]{results/spd_cpu_d=5.png}
        \caption{n=5}
        \label{fig:karcher_n5}
    \end{subfigure}
    \hfill
    \begin{subfigure}[b]{0.49\linewidth}
        \centering
        \includegraphics[width=\linewidth]{results/spd_cpu_d=10.png}
        \caption{n=10}
        \label{fig:karcher_n10}
    \end{subfigure}
    
    \caption{CPUtime result for Karcher mean problem.}
    \label{fig:Karcher_mean_unconstrained_cpu}
\end{figure}

\subsection{Additional Real Application Examples}
In this section, we provide additional numerical results to supplement the main paper. We present more examples of the DNN attack applications in Figures~\ref{fig:cifar_ex1}--\ref{fig:cifar_ex3}, and additional results for the Horizon leveling task in Figures~\ref{fig:horizon_ex1}--\ref{fig:horizon_ex3}.

It is important to note that the experimental settings, model architectures, and hyperparameters used in these additional experiments are identical to those described in the main paper.

% ==================================================================
% CIFAR EXAMPLE 1 (Cat -> Ship)
% ==================================================================
\begin{figure*}[htbp]
    \centering
    \setlength{\tabcolsep}{1pt}
    \renewcommand{\arraystretch}{0.5}

    % --- Part (a): Visual Images ---
    \begin{subfigure}[b]{\linewidth}
        \centering
        \begin{tabular}{ccccc}
            \includegraphics[width=0.15\linewidth]{results/CIFAR_result/9_Original_cat.png} &
            \includegraphics[width=0.15\linewidth]{results/CIFAR_result/9_PGD_ship.png} &
            \includegraphics[width=0.15\linewidth]{results/CIFAR_result/9_RGD_ship.png} &
            \includegraphics[width=0.15\linewidth]{results/CIFAR_result/9_RZO_ship.png} &
            \includegraphics[width=0.15\linewidth]{results/CIFAR_result/9_RDueling_ship.png} \\
            \tiny Original & \tiny PGD & \tiny RGD & \tiny ZO-RGD & \tiny RDNGD \\
        \end{tabular}
        \caption{Generated adversarial images}
        \label{fig:cifar_1_visual}
    \end{subfigure}
    
    \vspace{10pt} 

    \centering 
    
    \begin{subfigure}[b]{0.45\linewidth}
        \centering
        \includegraphics[width=\linewidth]{results/CIFAR_result/9_loss_vs_iteration.png}
        \caption{Adversarial Loss vs. Iteration}
        \label{fig:cifar_1_loss_iter}
    \end{subfigure}
    \hspace{0.5cm}
    \begin{subfigure}[b]{0.45\linewidth}
        \centering
        \includegraphics[width=\linewidth]{results/CIFAR_result/9_loss_vs_time.png}
        \caption{Adversarial Loss vs.  Time}
        \label{fig:cifar_1_loss_time}
    \end{subfigure}

    \caption{Additional attack result on CIFAR-10 (Example 1).}
    \label{fig:cifar_ex1}
\end{figure*}

% ==================================================================
% CIFAR EXAMPLE 2 (Ship -> Deer)
% ==================================================================
% ==================================================================
% CIFAR EXAMPLE 2 (Ship -> Deer)
% ==================================================================
\begin{figure*}[htbp]
    \centering
    \setlength{\tabcolsep}{1pt}
    \renewcommand{\arraystretch}{0.5}

    % --- Part (a): Visual Images ---
    \begin{subfigure}[b]{\linewidth}
        \centering
        \begin{tabular}{ccccc}
            \includegraphics[width=0.15\linewidth]{results/CIFAR_result/2_Original_ship.png} &
            \includegraphics[width=0.15\linewidth]{results/CIFAR_result/2_PGD_deer.png} &
            \includegraphics[width=0.15\linewidth]{results/CIFAR_result/2_RGD_ship.png} &
            \includegraphics[width=0.15\linewidth]{results/CIFAR_result/2_RZO_deer.png} &
            \includegraphics[width=0.15\linewidth]{results/CIFAR_result/2_RDueling_deer.png} \\
            \tiny Original & \tiny PGD & \tiny RGD & \tiny ZO-RGD & \tiny RDNGD \\
        \end{tabular}
        \caption{Generated adversarial images}
    \end{subfigure}

    \vspace{10pt}

    % --- Part (b) & (c): Loss Curves ---
    \centering
    \begin{subfigure}[b]{0.45\linewidth}
        \centering
        \includegraphics[width=\linewidth]{results/CIFAR_result/2_loss_vs_iteration.png}
        \caption{Adversarial Loss vs.  Iteration}
    \end{subfigure}
    \hspace{0.5cm}
    \begin{subfigure}[b]{0.45\linewidth}
        \centering
        \includegraphics[width=\linewidth]{results/CIFAR_result/2_loss_vs_time.png}
        \caption{Adversarial Loss vs.  Time}
    \end{subfigure}

    \caption{Additional attack result on CIFAR-10 (Example 2).}
    \label{fig:cifar_ex2}
\end{figure*}

% ==================================================================
% CIFAR EXAMPLE 3 (Ship -> Cat)
% ==================================================================
\begin{figure*}[htbp]
    \centering
    \setlength{\tabcolsep}{1pt}
    \renewcommand{\arraystretch}{0.5}

    % --- Part (a): Visual Images ---
    \begin{subfigure}[b]{\linewidth}
        \centering
        \begin{tabular}{ccccc}
            \includegraphics[width=0.15\linewidth]{results/CIFAR_result/3_Original_ship.png} &
            \includegraphics[width=0.15\linewidth]{results/CIFAR_result/3_PGD_cat.png} &
            \includegraphics[width=0.15\linewidth]{results/CIFAR_result/3_RGD_cat.png} &
            \includegraphics[width=0.15\linewidth]{results/CIFAR_result/3_RZO_cat.png} &
            \includegraphics[width=0.15\linewidth]{results/CIFAR_result/3_RDueling_cat.png} \\
            \tiny Original & \tiny PGD & \tiny RGD & \tiny ZO-RGD & \tiny RDNGD \\
        \end{tabular}
        \caption{Generated adversarial images}
    \end{subfigure}

    \vspace{10pt}

    % --- Part (b) & (c): Loss Curves ---
    \centering
    \begin{subfigure}[b]{0.45\linewidth}
        \centering
        \includegraphics[width=\linewidth]{results/CIFAR_result/3_loss_vs_iteration.png}
        \caption{Adversarial Loss vs.  Iteration}
    \end{subfigure}
    \hspace{0.5cm}
    \begin{subfigure}[b]{0.45\linewidth}
        \centering
        \includegraphics[width=\linewidth]{results/CIFAR_result/3_loss_vs_time.png}
        \caption{Adversarial Loss vs.  Time}
    \end{subfigure}

    \caption{Additional attack result on CIFAR-10 (Example 3).}
    \label{fig:cifar_ex3}
\end{figure*}

% ==================================================================
% HORIZON LEVELING (Single Figure with a,b,c cols)
% ==================================================================

% ==================================================================
% CIFAR PART (Keeping your previous settings)
% ==================================================================


% ==================================================================
% HORIZON LEVELING EXAMPLE 1
% ==================================================================
\begin{figure*}[htbp]
    \centering
    \begin{subfigure}[b]{0.32\linewidth}
        \centering
        \includegraphics[width=\linewidth]{results/Horizon_result/2600343901_0ff34b1428_o_original.png}
        \caption{Original Input}
    \end{subfigure}
    \hfill
    \begin{subfigure}[b]{0.32\linewidth}
        \centering
        \includegraphics[width=\linewidth]{results/Horizon_result/2600343901_0ff34b1428_o_correction.png}
        \caption{Corrected Result ($-9.17^{\circ}$)}
    \end{subfigure}
    \hfill
    \begin{subfigure}[b]{0.32\linewidth}
        \centering
        \includegraphics[width=\linewidth]{results/Horizon_result/2600343901_0ff34b1428_o_loss_vs_iteration.png}
        \caption{Loss vs Iteration}
    \end{subfigure}

    \caption{Additional horizon leveling result (Example 1).}
    \label{fig:horizon_ex1}
\end{figure*}

% ==================================================================
% HORIZON LEVELING EXAMPLE 2
% ==================================================================
\begin{figure*}[htbp]
    \centering
    \begin{subfigure}[b]{0.32\linewidth}
        \centering
        \includegraphics[width=\linewidth]{results/Horizon_result/2680766258_66b5a9fcbf_o_original.png}
        \caption{Original Input}
    \end{subfigure}
    \hfill
    \begin{subfigure}[b]{0.32\linewidth}
        \centering
        \includegraphics[width=\linewidth]{results/Horizon_result/2680766258_66b5a9fcbf_o_correction.png}
        \caption{Corrected Result ($20.05^{\circ}$)}
    \end{subfigure}
    \hfill
    \begin{subfigure}[b]{0.32\linewidth}
        \centering
        \includegraphics[width=\linewidth]{results/Horizon_result/2680766258_66b5a9fcbf_o_loss_vs_iteration.png}
        \caption{Loss vs Iteration}
    \end{subfigure}

    \caption{Additional horizon leveling result (Example 2).}
    \label{fig:horizon_ex2}
\end{figure*}

% ==================================================================
% HORIZON LEVELING EXAMPLE 3
% ==================================================================
\begin{figure*}[t]
    \centering
    \begin{subfigure}[b]{0.32\linewidth}
        \centering
        \includegraphics[width=\linewidth]{results/Horizon_result/1892424628_2bd1ef9423_o_original.png}
        \caption{Original Input}
    \end{subfigure}
    \hfill
    \begin{subfigure}[b]{0.32\linewidth}
        \centering
        \includegraphics[width=\linewidth]{results/Horizon_result/1892424628_2bd1ef9423_o_correction.png}
        \caption{Corrected Result ($-37.82^{\circ}$)}
    \end{subfigure}
    \hfill
    \begin{subfigure}[b]{0.32\linewidth}
        \centering
        \includegraphics[width=\linewidth]{results/Horizon_result/1892424628_2bd1ef9423_o_loss_vs_iteration.png}
        \caption{Loss vs Iteration}
    \end{subfigure}

    \caption{Additional horizon leveling result (Example 3).}
    \label{fig:horizon_ex3}
\end{figure*}

\end{document}



